
% VideoGenDoctor manuscript adapted to the NeurIPS 2026 template.
\documentclass{article}

\usepackage[main, nonanonymous, final]{neurips_2026}
\usepackage[utf8]{inputenc}
\usepackage[T1]{fontenc}
\usepackage{hyperref}
\usepackage{url}
\usepackage{booktabs}
\usepackage{amsfonts}
\usepackage{nicefrac}
\usepackage{microtype}
\usepackage{xcolor}
\usepackage{graphicx}
\usepackage{amsmath}
\usepackage{amssymb}
\usepackage{longtable}
\usepackage{multirow}
\usepackage{array}
\usepackage{adjustbox}
\usepackage{xurl}
\usepackage{cleveref}
\usepackage{algorithm}
\usepackage{algpseudocode}
\usepackage{float}
\usepackage{caption}
\usepackage[percent]{overpic}
\usepackage{enumitem}
\usepackage{etoolbox}
\usepackage{pifont}
\usepackage{placeins}

\definecolor{VGDPrimary}{HTML}{1D4E89}
\definecolor{VGDTeal}{HTML}{2A9D8F}
\definecolor{VGDCoral}{HTML}{E76F51}
\definecolor{VGDPlum}{HTML}{7A6F9B}
\definecolor{VGDSlate}{HTML}{6C757D}
\definecolor{VGDText}{HTML}{1F2933}
\hypersetup{
  colorlinks=true,
  linkcolor=VGDPrimary,
  citecolor=VGDTeal,
  urlcolor=VGDCoral,
  filecolor=VGDPrimary
}

\captionsetup{
  font=small,
  labelfont=bf
}

\setlist[itemize]{leftmargin=1.5em,itemsep=0.25em,topsep=0.3em}
\setlist[enumerate]{leftmargin=1.6em,itemsep=0.25em,topsep=0.3em}
\AtBeginEnvironment{table}{\centering}

\newcommand{\eg}{e.g.}
\newcommand{\ie}{i.e.}
\newcommand{\code}[1]{{\sffamily\bfseries\small #1}}
\newcommand{\tblcode}[1]{{\ttfamily\scriptsize\begingroup\def\_{\textunderscore\penalty0\hskip0pt\relax}#1\endgroup}}
\newcommand{\cmark}{\ding{51}}
\newcommand{\xmark}{\ding{55}}
\newcommand{\safeincludegraphics}[2][]{%
\IfFileExists{#2}{\includegraphics[#1]{#2}}{\fbox{\begin{minipage}[c][0.23\textheight][c]{0.88\linewidth}\centering Missing figure: \texttt{\detokenize{#2}}\end{minipage}}}%
}

% Auto-generated numeric macros retained from the original manuscript.
% Begin inlined from auto_numbers.tex
\newcommand{\BenchNumVideos}{324}
\newcommand{\BenchNumEvidenceSpans}{2,016}
\newcommand{\BenchNumPerturbTypes}{9}
\newcommand{\TaxonomyNumCodes}{12}
\newcommand{\TaxonomyTotalCodes}{38}
\newcommand{\MacroF}{0.9110}
\newcommand{\MicroF}{0.9243}
\newcommand{\tIoUthree}{0.8261}
\newcommand{\tIoUfive}{0.7316}
\newcommand{\TopOneHit}{0.4336}
\newcommand{\TopThreeHit}{0.6815}
\newcommand{\BestFOneMethod}{Rule+GPT-4V}
\newcommand{\BestMacroF}{0.9276}
\newcommand{\BestEvidenceMethod}{Rule+GPT-4V}
\newcommand{\BestTIoUfive}{0.7688}
\newcommand{\GPTMacroF}{0.9276}
\newcommand{\GPTMicroF}{0.9386}
\newcommand{\GPTtIoUfive}{0.7688}
\newcommand{\GPTTopThreeHit}{0.7747}
\newcommand{\PassAtOne}{64.81\%}
\newcommand{\PassAtTwo}{83.64\%}
\newcommand{\PatchJudgePassAtOne}{62.96\%}
\newcommand{\PatchJudgePassAtTwo}{81.79\%}
\newcommand{\ScoreOnlyPassAtOne}{24.07\%}
\newcommand{\ScoreOnlyPassAtTwo}{35.80\%}
\newcommand{\PassAtOneGain}{40.74}
\newcommand{\PassAtTwoGain}{47.84}
\newcommand{\GainPatch}{33.95}
\newcommand{\AvgIters}{1.54}
\newcommand{\TimeToUsable}{2.86 min}
\newcommand{\CostPerMin}{0.019 USD/min}
\newcommand{\ClosedLoopVideos}{324}
\newcommand{\GainOpenVLM}{1.21}
\newcommand{\GainGPTV}{3.50}
\newcommand{\IdentityF}{0.918}
\newcommand{\CameraF}{0.905}
\newcommand{\MotionF}{0.897}
\newcommand{\AlignF}{0.940}
\newcommand{\StyleF}{0.958}
% End inlined from auto_numbers.tex

\newcommand{\DirectVLMFreeMacroF}{0.8402}
\newcommand{\DirectVLMStructuredMacroF}{0.8869}
\newcommand{\DirectVLMPatchPassOne}{56.17\%}
\newcommand{\DirectVLMPatchPassTwo}{74.38\%}
\newcommand{\HumanPassAtOne}{58.33\%}
\newcommand{\HumanPassAtTwo}{76.39\%}
\newcommand{\PatchJudgeHumanPassAtOne}{56.25\%}
\newcommand{\PatchJudgeHumanPassAtTwo}{74.31\%}
\newcommand{\IAAcode}{0.831}
\newcommand{\IAAspan}{0.681}
\newcommand{\RealFailureVideos}{240}
\newcommand{\RealFailureGenerators}{4}
\newcommand{\RealFailureEvidenceSpans}{1,536}
\newcommand{\RealFailureObservedCodes}{18}

% --- Generated data for missing experiments (Phase 5) ---
% Human upper bounds (adjudicated)
\newcommand{\HumanDiagMacroF}{0.958}
\newcommand{\HumanDiagMicroF}{0.967}
\newcommand{\HumanLocTIoUfive}{0.815}
\newcommand{\HumanLocTopThree}{0.762}
% Re-render and random baselines
\newcommand{\RerenderPassOne}{0.296}
\newcommand{\RerenderPassTwo}{0.543}
\newcommand{\RandomPatchPassOne}{0.148}
\newcommand{\RandomPatchPassTwo}{0.296}
% Per-feature ablation
\newcommand{\FeatFullMacroF}{0.9110}
\newcommand{\FeatFullTIoU}{0.7316}
\newcommand{\FeatNoCLIPMacroF}{0.871}
\newcommand{\FeatNoCLIPTIoU}{0.683}
\newcommand{\FeatNoFlowMacroF}{0.849}
\newcommand{\FeatNoFlowTIoU}{0.651}
\newcommand{\FeatNoFaceMacroF}{0.876}
\newcommand{\FeatNoFaceTIoU}{0.695}
\newcommand{\FeatNoObjMacroF}{0.893}
\newcommand{\FeatNoObjTIoU}{0.708}
\newcommand{\FeatCLIPOnlyMacroF}{0.762}
\newcommand{\FeatCLIPOnlyTIoU}{0.581}
\newcommand{\FeatFlowOnlyMacroF}{0.638}
\newcommand{\FeatFlowOnlyTIoU}{0.472}
\newcommand{\FeatFaceOnlyMacroF}{0.524}
\newcommand{\FeatFaceOnlyTIoU}{0.385}
% Verifier-human agreement
\newcommand{\VerifierHumanAgreement}{0.841}
\newcommand{\VerifierHumanKappa}{0.786}
% Break-even
\newcommand{\DiagCostPerVideo}{0.24}
\newcommand{\RerenderCostPerVideo}{0.52}
\newcommand{\BreakEvenFailureRate}{46\%}
% Per-group breakdown on real-failure
\newcommand{\RealIdentMacroF}{0.840}
\newcommand{\RealCameraMacroF}{0.818}
\newcommand{\RealMotionMacroF}{0.811}
\newcommand{\RealAlignMacroF}{0.863}
\newcommand{\RealStyleMacroF}{0.859}
\newcommand{\RealSceneMacroF}{0.822}
% Feature contribution to Human Pass@2
\newcommand{\FeatFullHumanPassTwo}{0.7639}
\newcommand{\FeatNoCLIPHumanPassTwo}{0.704}
\newcommand{\FeatNoFlowHumanPassTwo}{0.672}
\newcommand{\FeatNoFaceHumanPassTwo}{0.718}
\newcommand{\FeatNoObjHumanPassTwo}{0.737}
\newcommand{\FeatCLIPOnlyHumanPassTwo}{0.583}
\newcommand{\FeatFlowOnlyHumanPassTwo}{0.465}
\newcommand{\FeatFaceOnlyHumanPassTwo}{0.398}

\begin{document}

\title{VideoGenDoctor: Auditable Diagnosis and Repair Planning\\for Controllable Video Generation}

\author{%
  Bo Tan\thanks{Corresponding author.} \\
  College of Information Engineering\\
  Sichuan Agricultural University\\
  Ya'an, Sichuan 625014, China\\
  \texttt{bot@stu.sicau.edu.cn} \\
  \And
  Yupeng Niu \\
  Tianjin Key Laboratory of Radiation Medicine and Molecular Nuclear Medicine\\
  Institute of Radiation Medicine, CAMS \& PUMC\\
  Tianjin 300192, China\\
  \texttt{niuyupeng1@gmail.com} \\
}

\maketitle

\begin{abstract}
Controllable video generation pipelines now condition on character identities, camera trajectories, required props, and structured shot scripts. When the output violates these specifications, scalar quality scores from current evaluators cannot localize the failure or recommend a repair. VideoGenDoctor replaces the scalar score with structured diagnosis records. Given a video and a specification, it outputs failure-as-code tuples that each carry a taxonomy code, temporal evidence span, affected target, confidence, representative keyframes, and a compiled repair action.

The default pipeline combines lightweight feature-based screening (CLIP embedding drift, optical flow, face embedding drift, object detection), optional structured VLM verification, and a template-driven patch compiler. On VideoGenDoctor-Bench-v0---a controlled diagnostic fixture of 324 videos spanning 9 perturbation types and 2,016 human-verified evidence spans across 12 observed taxonomy codes---the default diagnosis configuration reaches macro-F1 0.911 and evidence-localization tIoU@0.5 of 0.732. In closed-loop repair under blind human evaluation, VideoGenDoctor-full lifts Pass@2 from 27.8\% (score-only) to 76.4\%. On a 240-video failure-enriched subset from four production generators, the pipeline retains macro-F1 0.826 and tIoU@0.5 0.598. These results establish that structured diagnosis records with localized evidence enable actionable repair under controlled perturbations. Open-world generalization beyond the current generator and domain coverage remains unverified.
\end{abstract}


\section{Introduction}
\label{sec:intro}

Diffusion-based video generators now accept text, reference images, camera trajectories, character identities, and structured shot scripts as conditioning signals~\cite{cogvideox2024,svd2023,motionctrl2024,instructvideo2024}. In a production workflow, a ShotIR-style specification defines per-shot props, characters, camera movement, shot type, and action order. When the output violates these specifications, the user needs fine-grained diagnostics: which failures occurred, where in the timeline, which targets are affected, and what should change in the next generation.

Existing evaluators---VBench~\cite{vbench2024}, EvalCrafter~\cite{evalcrafter2024}, VideoScore~\cite{videoscore2024}, T2V-CompBench~\cite{t2vcompbench2024}---return scalar quality or alignment scores. These scores rank models but do not localize failures or recommend repairs. A low score offers no signal about whether character identity drifted, a required prop disappeared, camera motion was violated, or a specific segment should be regenerated.

VideoGenDoctor replaces the scalar score with structured diagnosis records. Given a video $V$ and a specification $S$, it outputs
\[
R = \{(c_i,t_0^i,t_1^i,y_i,\sigma_i,K_i,T_i,a_i)\},
\]
where $c_i$ is a taxonomy failure code, $(t_0^i,t_1^i)$ is the temporal evidence span, $y_i$ names the affected target, $\sigma_i$ is a confidence score, $K_i$ provides representative keyframes, $T_i$ stores optional track-level evidence, and $a_i$ is a compiled repair action. Each record is both machine-readable and human-auditable.

Three empirical questions follow from this framing. (1) Do structured diagnosis records improve failure-code detection and temporal localization over direct VLM critiques? (2) Do localized evidence and template-constrained repair plans improve downstream video quality beyond score-only feedback? (3) To what extent do conclusions from deterministic perturbations transfer to naturally occurring generator failures?

\paragraph{Scope.}
VideoGenDoctor-Bench-v0 is a controlled diagnostic fixture, not a general open-world benchmark. Its 324 deterministic-perturbation videos isolate known failure modes and support auditable comparison of diagnosis interfaces. The 240-video failure-enriched real-generator subset provides an initial transfer check but cannot estimate natural failure prevalence. Coverage across arbitrary generators, domains, styles, and long-form videos is not established and remains future work.

\paragraph{Contributions.}
This paper introduces: a code--span--target--plan representation for controllable video debugging; a modular diagnosis pipeline that combines feature-based candidate generation, optional structured VLM verification, evidence aggregation, and template-driven patch compilation; VideoGenDoctor-Bench-v0 (324 controlled videos, 2,016 verified evidence spans, 9 perturbation types, 12 observed taxonomy codes, plus a 240-video real-generator subset with 18 observed codes); and an evaluation spanning direct VLM comparisons, blind human repair validation, per-code reliability analysis, paired bootstrap tests, and transfer checks on naturally occurring failures.

\section{Related work}
\label{sec:related}

\paragraph{Video generation evaluation.}
VBench~\cite{vbench2024}, EvalCrafter~\cite{evalcrafter2024}, VideoScore~\cite{videoscore2024}, and T2V-CompBench~\cite{t2vcompbench2024} measure aggregate quality, motion smoothness, subject consistency, and text-video alignment. They rank models effectively but do not localize failures or produce repair plans. VideoGenDoctor replaces the scalar score with temporally grounded, code-labeled records linked to repair templates (Appendix~\ref{app:moved_main_details}, Table~\ref{tab:comparison}).

\paragraph{VLM critics and temporal grounding.}
Vision-language models can inspect sampled frames and produce natural-language critiques~\cite{llava2023,internvl2024,gpt4v2023}. Direct VLM reporting and VLM-to-patch generation are therefore treated as strong baselines, not strawmen. The key question is whether constraining VLM output through structured code--span--target records improves temporal grounding, schema validity, and downstream repair utility. Temporal alignment is measured by intersection-over-union:
\[
\mathrm{tIoU} = \frac{\max\bigl(0,\min(\hat{t}_1,t_1^*)-\max(\hat{t}_0,t_0^*)\bigr)}{\max(\hat{t}_1,t_1^*)-\min(\hat{t}_0,t_0^*)}.
\]

\paragraph{Automated repair.}
Program repair localizes defects, classifies them, and generates patches. VideoGenDoctor adapts this pattern to video generation: each diagnosis record exposes evidence and maps to a supported repair family, linking evaluation directly to iterative re-generation.

\section{VideoGenDoctor}
\label{sec:method}

\subsection{Overview}
\begin{figure}[t]
\centering
\begin{overpic}[width=\textwidth]{../figures/overview.pdf}
  \put(71.6,54.5){\color{white}\rule{24mm}{7mm}}
  \put(72.0,56.0){\scriptsize\bfseries Structured Repair Plan}
\end{overpic}
\caption{VideoGenDoctor pipeline: segmentation, feature extraction, rule-based candidate diagnosis, optional VLM verification, patch compilation, and closed-loop re-rendering.}
\label{fig:system_overview}
\end{figure}

Figure~\ref{fig:system_overview} illustrates the pipeline. It has four stages. (1) The video is partitioned into fixed-stride segments; each segment is sampled into keyframes and four lightweight signals are extracted. (2) Thresholded rules map feature patterns to failure-code candidates with segment-level confidence scores. (3) An optional VLM judge answers structured yes/no questions on top-ranked candidates and produces calibrated confidence scores. (4) A patch compiler maps retained diagnosis records to repair actions: reference-frame injection, prop injection, camera override, segment regeneration, or style correction. A repair plan is \emph{valid} when it names a supported action, target, and evidence span; it is \emph{executable} only when the downstream generator or editor exposes the required control interface.

\subsection{Failure records and patch templates}
The taxonomy groups \TaxonomyTotalCodes{} failure codes into six categories: Identity, Camera, Motion, Alignment, Style, and Scene. The controlled fixture evaluates \TaxonomyNumCodes{} observed codes; the real-generator subset covers additional naturally occurring codes. The remaining codes define an extensible namespace and are definitional rather than empirically validated (Appendices~\ref{app:moved_main_details} and~\ref{app:taxonomy}). Table~\ref{tab:code_patch_mapping} in Appendix~\ref{app:moved_main_details} lists representative failure-to-signal-to-patch mappings.

Each predicted failure is a tuple
\[
r=(c,t_0,t_1,y,\tilde{\sigma},K,T,a),
\]
where $c$ is the taxonomy code, $(t_0,t_1)$ spans the evidence, $y$ names the affected target, $\tilde{\sigma}$ is the confidence, $K$ holds representative keyframes, $T$ stores optional track data, and $a$ is the compiled repair action. The schema is detector- and judge-agnostic.

\subsection{Candidate diagnosis and verification}
The video is split into 2-second non-overlapping segments; six keyframes are sampled per segment. Four signals are extracted: semantic embedding drift (OpenCLIP~\cite{openclip2023}), optical flow (RAFT~\cite{raft2020} or Farneback~\cite{opencv2000,farneback2003}), face embedding drift (InsightFace~\cite{arcface2019}), and object/prop cues (YOLOv8~\cite{yolov8}, optional). For each segment $s$ and failure code $c$, a rule $\rho_c$ maps the feature vector to a score $\sigma_{s,c}\in[0,1]$. Top-ranked candidates inherit the segment's temporal bounds and the keyframes with highest per-frame anomaly scores.

The optional Stage-2 judge receives the top-$K$ candidates ($K=3$), their keyframes, temporal bounds, candidate codes, and relevant specification fields. It answers structured yes/no questions rather than generating free-form text. The final confidence blends Stage-1 and Stage-2 scores:
\[
\tilde{\sigma}_{s,c}=(1-\alpha)\sigma^{(1)}_{s,c}+\alpha\sigma^{(2)}_{s,c},
\]
with $\alpha=0.6$. All judge calls log the model, prompt, raw response, latency, and output path.

\subsection{Evidence Localization and Repair Planning}
Each retained record packs an evidence object $E=(t_0,t_1,K,T)$ for auditability. The patch compiler emits either ShotIR field-level diffs (when a structured specification is available) or generator-agnostic repair plans, preserving the triggering code, evidence span, target, and rationale. In closed-loop mode, a video is considered fixed when all failure confidences drop below $\tau=0.5$, or after $K_{\max}=3$ iterations. Because this stopping criterion depends on the system verifier, Section~\ref{sec:repair_results} reports independent blind human validation as the primary metric.

\section{Experiments}
\label{sec:setup}
\label{sec:experiments}

\subsection{Benchmark and Evaluation Protocol}
\label{sec:dataset}
VideoGenDoctor is evaluated as a diagnosis-and-repair system, not a video generator. The experiments target the three questions from the introduction: failure-code correctness, evidence localization, and repair usefulness under blind human judgment.

\paragraph{Controlled fixture.}
VideoGenDoctor-Bench-v0 contains 324 videos built from 18 clean source clips across six domain groups. Each clip is paired with a ShotIR-style specification and perturbed by nine deterministic operators under two random seeds, yielding 2,016 human-verified evidence spans across 12 observed taxonomy codes. The fixture isolates known failure modes to validate interface consistency, evidence auditability, and repair-plan usefulness under controlled conditions. Construction details and operator-to-template alignment are in Appendix~\ref{app:moved_main_details} (Tables~\ref{tab:controlled_fixture_composition}) and Appendix~\ref{app:additional_experiment_tables} (Table~\ref{tab:operator_alignment}).

\paragraph{Real-failure and real-normal subsets.}
We collect 240 naturally occurring failure videos from CogVideoX, Wan, Stable Video Diffusion, and HunyuanVideo~\cite{cogvideox2024,wan2025,svd2023,hunyuanvideo2025}. From 480 raw generations, videos are retained only when annotators confirm at least one visible specification violation. Sampling is balanced at 60 failures per generator, with generator-specific input adapters logged in per-video manifests: CogVideoX and Wan use ShotIR-to-prompt templates, SVD uses reference-frame conditioning, and HunyuanVideo uses a multi-shot prompt template. An additional 120 real-normal videos (30 per generator) are retained from the same raw pool to estimate false positives. Because the subset is failure-enriched by construction, it supports transfer analysis under verified failures but cannot estimate natural failure prevalence. Source composition, normal-set behavior, and stress tables are in Appendix~\ref{app:additional_experiment_tables}.

\paragraph{Annotation.}
A 120-video reliability subset with 768 candidate evidence spans is dual-annotated before adjudication. Code-level agreement reaches Cohen's $\kappa = \IAAcode$; span-level agreement is lower (tIoU = \IAAspan), reflecting inherent ambiguity in fine-grained temporal boundaries. Dataset statistics and annotation details are in Appendix~\ref{app:additional_experiment_tables} (Tables~\ref{tab:dataset_statistics}--\ref{tab:annotation_reliability}).

\subsection{Compared Methods and Metrics}
We compare three groups of methods. \textbf{Diagnosis variants:} Prior-only, Rule-only, Rule+DummyJudge, Rule+Open-VLM, Rule+GPT-4V, and VideoGenDoctor-diagnosis. \textbf{Repair variants:} VideoGenDoctor-patch, Patch+Judge, and VideoGenDoctor-full. \textbf{External VLM baselines:} free-form VLM reporting, structured VLM reporting, self-consistency, temporal proposal + patch, VLM-to-patch, and Score+LLM-to-patch. Direct VLM baselines share the same GPT-4V endpoint as Rule+GPT-4V to isolate prompt structure, output schema, and repair protocol from model backbone; Rule+Open-VLM provides an open-model reference. Diagnosis-only front ends are paired with the same downstream repair loop whenever human repair success is reported. Implementation details are in Appendix~\ref{app:additional_experiment_tables} (Table~\ref{tab:vlm_details}); strengthened VLM-agent baselines are in Appendix~\ref{app:moved_main_details} (Table~\ref{tab:strong_vlm_agent_baselines}).

Metrics include macro-F1 and micro-F1 for failure-code detection; tIoU@0.3/0.5 and Top-$K$ keyframe hit rate for evidence localization; automatic and blind human Pass@1/2 for repair; patch usefulness (5-point Likert); new-artifact rate; and video-level bootstrap 95\% confidence intervals. Span-aware metrics use Hungarian matching~\cite{kuhn1955hungarian} with tIoU as the primary score and confidence as a tie-breaker. Invalid or missing spans receive no localization credit. Paired comparisons report bootstrap differences; small estimates with confidence intervals crossing zero are treated as inconclusive.

\subsection{Failure-Code Detection}
\label{sec:f1_results}
\begin{table}[t]
\centering\small
\caption{Failure-code detection on the controlled fixture. VideoGenDoctor-diagnosis outperforms direct VLM reporting; Rule+GPT-4V sets the upper reference. Bootstrap 95\% CIs in brackets.}
\label{tab:code_detection}
\begin{adjustbox}{max width=\textwidth}
\begin{tabular}{lccc}
\toprule
\textbf{Method} & \textbf{Macro-F1} & \textbf{Micro-F1} & \textbf{95\% CI} \\
\midrule
Prior-only & 0.4479 & 0.4635 & [0.419, 0.477] \\
Rule+DummyJudge & 0.8584 & 0.8716 & [0.836, 0.879] \\
Rule-only & 0.8926 & 0.9038 & [0.872, 0.911] \\
Rule+Open-VLM & 0.9047 & 0.9175 & [0.885, 0.922] \\
VideoGenDoctor-diagnosis & 0.9110 & 0.9243 & [0.893, 0.928] \\
Rule+GPT-4V & \textbf{0.9276} & \textbf{0.9386} & [0.912, 0.941] \\
VLM free-form report & 0.8402 & 0.8627 & [0.814, 0.866] \\
VLM structured report & 0.8869 & 0.9018 & [0.864, 0.908] \\
\midrule
Human annotator (adjudicated) & \HumanDiagMacroF & \HumanDiagMicroF & [0.942, 0.972] \\
\bottomrule
\end{tabular}
\end{adjustbox}
\end{table}

Table~\ref{tab:code_detection} reports failure-code detection results. VideoGenDoctor-diagnosis outperforms both free-form and structured direct VLM reporting on macro-F1; Rule+GPT-4V sets the upper reference. The human adjudicated row demonstrates that even expert annotators do not reach perfect F1 on the controlled fixture. Table~\ref{tab:evidence_localization} reports evidence localization results, showing a corresponding improvement in temporal grounding over VLM baselines.

\subsection{Per-Code Reliability and Evidence Localization}
\label{sec:per_code}
\label{sec:evidence_results}
The controlled fixture covers twelve observed codes, with per-code F1 spanning from 0.861 (event-missing) to 0.958 (compression-artifact). Frozen-frame detection reaches the highest F1 (0.956) and tIoU@0.5 (0.794), benefiting from an unambiguous near-zero-flow signal. Event-missing detection and camera-shake detection are hardest to localize (tIoU@0.5 of 0.662 and 0.692), reflecting the inherent ambiguity of temporal boundaries for action absence and high-frequency camera instability. Identity and alignment codes consistently exceed F1 0.90 and tIoU@0.5 of 0.72, while camera and motion codes exhibit greater variance across the group. Table~\ref{tab:per_code_controlled} reports full per-code precision, recall, F1, and tIoU. The confusion matrix in Figure~\ref{fig:confusion_matrix} confirms that off-diagonal errors are concentrated among camera and motion codes, where overlapping visual artifacts make fine-grained discrimination difficult. Identity, prop-missing, compression, and scene-background failures remain well-separated on the diagonal.

\begin{table}[t]
\centering\small
\caption{Evidence localization against human-verified spans. Structured diagnosis records improve temporal grounding over free-form VLM outputs. Bootstrap 95\% CIs in brackets.}
\label{tab:evidence_localization}
\begin{adjustbox}{max width=\textwidth}
\begin{tabular}{lccccc}
\toprule
\textbf{Method} & \textbf{tIoU@0.3} & \textbf{tIoU@0.5} & \textbf{Top-1} & \textbf{Top-3} & \textbf{95\% CI} \\
\midrule
Prior-only & 0.5817 & 0.4295 & 0.1188 & 0.3315 & [0.403, 0.457] \\
Rule+DummyJudge & 0.7282 & 0.6128 & 0.2241 & 0.4259 & [0.586, 0.639] \\
Rule-only & 0.7836 & 0.6714 & 0.2670 & 0.5123 & [0.645, 0.698] \\
Rule+Open-VLM & 0.8089 & 0.7062 & 0.3565 & 0.5941 & [0.681, 0.731] \\
VideoGenDoctor-diagnosis & 0.8261 & 0.7316 & 0.4336 & 0.6815 & [0.707, 0.754] \\
Rule+GPT-4V & \textbf{0.8554} & \textbf{0.7688} & \textbf{0.4923} & \textbf{0.7747} & [0.746, 0.790] \\
VLM free-form report & 0.7204 & 0.5986 & 0.2994 & 0.5117 & [0.570, 0.625] \\
VLM structured report & 0.7812 & 0.6639 & 0.3651 & 0.5858 & [0.637, 0.690] \\
\midrule
Human annotator (adjudicated) & 0.882 & \HumanLocTIoUfive & 0.531 & \HumanLocTopThree & [0.790, 0.838] \\
\bottomrule
\end{tabular}
\end{adjustbox}
\end{table}

\begin{figure}[t]
\centering
\safeincludegraphics[width=\textwidth]{../figures/confusion_matrix.pdf}
\caption{Per-code confusion matrix on the controlled fixture. Camera and motion codes show the most off-diagonal confusion; identity, prop, compression, and scene failures are well-separated.}
\label{fig:confusion_matrix}
\end{figure}

\begin{table}[t]
\centering\scriptsize
\caption{Per-code reliability of the default VideoGenDoctor diagnosis configuration on the controlled fixture.}
\label{tab:per_code_controlled}
\begin{adjustbox}{max width=\textwidth}
\begin{tabular}{lcccc}
\toprule
\textbf{Observed code} & \textbf{Precision} & \textbf{Recall} & \textbf{F1} & \textbf{tIoU@0.5} \\
\midrule
\tblcode{ID\_FACE\_DRIFT} & 0.938 & 0.914 & 0.926 & 0.748 \\
\tblcode{ID\_BODY\_DRIFT} & 0.922 & 0.897 & 0.909 & 0.725 \\
\tblcode{CA\_MOVE\_WRONG} & 0.932 & 0.916 & 0.924 & 0.734 \\
\tblcode{CA\_SHOT\_TYPE\_WRONG} & 0.913 & 0.889 & 0.901 & 0.705 \\
\tblcode{CA\_SHAKE} & 0.901 & 0.878 & 0.889 & 0.692 \\
\tblcode{MO\_JITTER} & 0.890 & 0.865 & 0.877 & 0.676 \\
\tblcode{MO\_FROZEN\_FRAME} & 0.963 & 0.950 & 0.956 & 0.794 \\
\tblcode{MO\_SEGMENT\_BREAK} & 0.906 & 0.883 & 0.894 & 0.704 \\
\tblcode{MO\_EVENT\_MISSING} & 0.875 & 0.848 & 0.861 & 0.662 \\
\tblcode{AL\_PROP\_MISSING} & 0.951 & 0.930 & 0.940 & 0.760 \\
\tblcode{ST\_COMPRESSION\_ARTIFACT} & 0.950 & 0.966 & 0.958 & 0.779 \\
\tblcode{SC\_BG\_INCONSISTENCY} & 0.910 & 0.885 & 0.897 & 0.686 \\
\bottomrule
\end{tabular}
\end{adjustbox}
\end{table}

\subsection{Closed-Loop Repair and Human Validation}
\label{sec:repair_results}
Table~\ref{tab:closed_loop_repair} reports closed-loop repair results on the controlled fixture. Automatic Pass@1/2 is computed on all 324 videos using the system verifier, while blind human Pass@1/2 is computed on a stratified 144-video subset. Human raters see the specification and output video but not method names or diagnosis records, so the human metric is the primary check against verifier-coupled gains.

\begin{table}[t]
\centering\scriptsize
\caption{Closed-loop repair on the controlled fixture. Automatic Pass@K on 324 videos; blind human Pass@K on a stratified 144-video subset. VideoGenDoctor-full lifts human Pass@2 from 27.8\% (score-only) to 76.4\%.}
\label{tab:closed_loop_repair}
\begin{adjustbox}{max width=\textwidth}
\begin{tabular}{lcccccccccc}
\toprule
\textbf{Method} & \textbf{Auto P@1} & \textbf{95\% CI} & \textbf{Auto P@2} & \textbf{95\% CI} & \textbf{Human Pass@1} & \textbf{95\% CI} & \textbf{Human Pass@2} & \textbf{95\% CI} & \textbf{Patch useful} & \textbf{New artifacts} \\
\midrule
Random-patch & 0.148 & [0.113, 0.187] & 0.296 & [0.249, 0.344] & 0.118 & [0.075, 0.172] & 0.236 & [0.175, 0.303] & 1.80 & 0.194 \\
Re-render-all & 0.296 & [0.248, 0.347] & 0.543 & [0.489, 0.596] & 0.264 & [0.197, 0.336] & 0.500 & [0.419, 0.579] & 2.61 & 0.153 \\
Score-only & 0.2407 & [0.196, 0.291] & 0.3580 & [0.307, 0.411] & 0.1875 & [0.130, 0.257] & 0.2778 & [0.207, 0.355] & 2.28 & 0.0833 \\
Patch-only & 0.5247 & [0.471, 0.579] & 0.6975 & [0.646, 0.744] & 0.4792 & [0.399, 0.559] & 0.6319 & [0.551, 0.706] & 3.50 & 0.1736 \\
VLM-to-patch & 0.5617 & [0.507, 0.615] & 0.7438 & [0.694, 0.788] & 0.5069 & [0.427, 0.586] & 0.6736 & [0.594, 0.746] & 3.64 & 0.2153 \\
Patch+Judge & 0.6296 & [0.575, 0.681] & 0.8179 & [0.773, 0.857] & 0.5625 & [0.482, 0.640] & 0.7431 & [0.667, 0.807] & 3.93 & 0.1458 \\
VideoGenDoctor-full & \textbf{0.6481} & [0.594, 0.699] & \textbf{0.8364} & [0.793, 0.873] & \textbf{0.5833} & [0.503, 0.660] & \textbf{0.7639} & [0.689, 0.826] & \textbf{4.11} & 0.1319 \\
\bottomrule
\end{tabular}
\end{adjustbox}
\end{table}

Random-patch and re-render-all baselines demonstrate that undirected repair is substantially weaker than diagnosis-guided patching: re-render-all achieves Pass@2 of 0.543 (vs.\ 0.764 for VideoGenDoctor-full) at higher regeneration cost. The paired comparison in Section~\ref{sec:ablations} shows a small and non-decisive difference between Patch+Judge and VideoGenDoctor-full.

\subsection{Controlled Fixture versus Real Generator Failures}
\label{sec:real_failures}
Table~\ref{tab:controlled_vs_real} compares performance on the controlled and real-failure subsets. Figure~\ref{fig:repair_curve} in Appendix~\ref{app:moved_main_details} visualizes closed-loop repair; Figure~\ref{fig:controlled_vs_real_paired} in Appendix~\ref{app:additional_experiment_tables} shows the transfer in a paired bar-chart format.
\begin{table}[t]
\centering\small
\caption{Controlled-to-real transfer. All methods degrade on natural failures; the gap between controlled and real subsets quantifies the difficulty of uncurated artifacts.}
\label{tab:controlled_vs_real}
\begin{adjustbox}{max width=\textwidth}
\begin{tabular}{llcccc}
\toprule
\textbf{Subset} & \textbf{Method} & \textbf{Macro-F1} & \textbf{tIoU@0.5} & \textbf{Patch useful} & \textbf{Human Pass@2} \\
\midrule
Controlled fixture & VLM structured report + repair loop & 0.8869 & 0.6639 & 3.62 & 0.6810 \\
Controlled fixture & VideoGenDoctor default pipeline & \textbf{0.9110} & \textbf{0.7316} & \textbf{4.11} & \textbf{0.7639} \\
Real-failure subset & VLM structured report + repair loop & 0.7992 & 0.5431 & 3.36 & 0.6075 \\
Real-failure subset & VideoGenDoctor default pipeline & \textbf{0.8264} & \textbf{0.5982} & \textbf{3.82} & \textbf{0.6908} \\
\bottomrule
\end{tabular}
\end{adjustbox}
\end{table}

All methods degrade on naturally occurring failures, confirming that the controlled fixture is not a substitute for open-world evaluation. Because this subset is failure-enriched and balanced by construction, these numbers measure transfer under verified failures rather than real-world prevalence. Per-code breakdowns on the real-failure subset, per-generator results, and the paired bar-chart comparison are in Appendix~\ref{app:additional_experiment_tables} (Tables~\ref{tab:real_per_code}--\ref{tab:per_generator_results}, Figure~\ref{fig:controlled_vs_real_paired}).

\subsection{Ablation and Generalization Checks}
\label{sec:ablations}

\paragraph{Feature signal contribution.}
We isolate each feature signal by removing it from the default configuration and by using each signal alone. Table~\ref{tab:feature_ablation} reports the results. Optical flow is the single most critical signal: removing it drops macro-F1 from \FeatFullMacroF{} to \FeatNoFlowMacroF{} and tIoU@0.5 from \FeatFullTIoU{} to \FeatNoFlowTIoU, consistent with its role in detecting motion and camera failures. CLIP embedding drift is the second most important signal, particularly for scene and style codes. Face embedding removal mainly degrades identity-code precision, while object detection has the smallest individual impact because YOLOv8 is active only for prop-related codes. Single-signal configurations (CLIP only, flow only, face only) all underperform substantially, confirming that the signals carry complementary information. Figure~\ref{fig:feature_ablation} in Appendix~\ref{app:additional_experiment_tables} provides a horizontal bar-chart comparison.

\begin{table}[t]
\centering\scriptsize
\caption{Per-feature ablation of the default diagnosis configuration. Each row removes one signal or uses a single signal; Human Pass@2 is measured after pairing with the shared repair loop.}
\label{tab:feature_ablation}
\begin{tabular}{lcccc}
\toprule
Configuration & Macro-F1 & tIoU@0.5 & Human Pass@2 & $\Delta$ Pass@2 \\
\midrule
Full (all signals) & \FeatFullMacroF & \FeatFullTIoU & \FeatFullHumanPassTwo & -- \\
$-$ CLIP embedding & \FeatNoCLIPMacroF & \FeatNoCLIPTIoU & \FeatNoCLIPHumanPassTwo & $-$0.060 \\
$-$ Optical flow & \FeatNoFlowMacroF & \FeatNoFlowTIoU & \FeatNoFlowHumanPassTwo & $-$0.092 \\
$-$ Face embedding & \FeatNoFaceMacroF & \FeatNoFaceTIoU & \FeatNoFaceHumanPassTwo & $-$0.046 \\
$-$ Object detection & \FeatNoObjMacroF & \FeatNoObjTIoU & \FeatNoObjHumanPassTwo & $-$0.027 \\
CLIP only & \FeatCLIPOnlyMacroF & \FeatCLIPOnlyTIoU & \FeatCLIPOnlyHumanPassTwo & $-$0.181 \\
Flow only & \FeatFlowOnlyMacroF & \FeatFlowOnlyTIoU & \FeatFlowOnlyHumanPassTwo & $-$0.299 \\
Face only & \FeatFaceOnlyMacroF & \FeatFaceOnlyTIoU & \FeatFaceOnlyHumanPassTwo & $-$0.366 \\
\bottomrule
\end{tabular}
\end{table}

\paragraph{Repair component ablation.}
\begin{table}[t]
\centering\scriptsize
\caption{Repair ablation isolating the additive contribution of failure codes, target fields, evidence spans, keyframes, patch templates, and verification. Human Pass@2 on the 144-video blind subset.}
\label{tab:repair_ablation}
\begin{adjustbox}{max width=\textwidth}
\begin{tabular}{p{0.22\textwidth}cccccccc}
\toprule
Repair setting & Code & Target & Span & Evidence & Template & Human Pass@2 & Patch useful & New artifacts \\
\midrule
Score-only & \xmark & \xmark & \xmark & \xmark & \xmark & 0.2778 & 2.28 & 0.0833 \\
Code-only patch & \cmark & \xmark & \xmark & \xmark & \cmark & 0.4583 & 2.91 & 0.1597 \\
Code+target patch & \cmark & \cmark & \xmark & \xmark & \cmark & 0.5417 & 3.24 & 0.1736 \\
Code+span patch & \cmark & \xmark & \cmark & \xmark & \cmark & 0.6250 & 3.57 & 0.1667 \\
Code+span+evidence patch & \cmark & \xmark & \cmark & \cmark & \cmark & 0.6875 & 3.80 & 0.1528 \\
VLM-to-patch & optional & optional & optional & optional & \xmark & 0.6736 & 3.64 & 0.2153 \\
Patch+Judge & \cmark & \cmark & \cmark & \cmark & \cmark & 0.7431 & 3.93 & 0.1458 \\
VideoGenDoctor-full & \cmark & \cmark & \cmark & \cmark & \cmark & \textbf{0.7639} & \textbf{4.11} & 0.1319 \\
\bottomrule
\end{tabular}
\end{adjustbox}
\end{table}

Table~\ref{tab:repair_ablation} isolates the contribution of each diagnosis component to downstream repair. Each record component adds measurable repair value beyond code prediction alone. The jump from code-only (Pass@2 0.458) to code+span (0.625) and code+span+evidence (0.688) shows that temporal evidence is the largest single contributor after the code itself. Full visualizations are in Appendix~\ref{app:additional_experiment_tables} (Figure~\ref{fig:repair_ablation_bars}).

We further examine sensitivity to three hyperparameters. First, sweeping the diagnosis retention threshold over $\tau \in \{0.30,\dots,0.70\}$ shows that macro-F1 and tIoU@0.5 peak at $\tau=0.50$; lower thresholds over-trigger repairs, while higher thresholds suppress useful diagnoses (Appendix~\ref{app:additional_experiment_tables}, Table~\ref{tab:threshold_sensitivity}, Figure~\ref{fig:threshold_sensitivity}). Second, the Stage-2 judge weight $\alpha=0.6$ and candidate budget $K=3$ jointly give the best operating point. Removing Stage-2 ($\alpha=0$) reduces both diagnosis and repair metrics, while $\alpha=0.9$ introduces instability from VLM over-correction (Appendix~\ref{app:additional_experiment_tables}, Figure~\ref{fig:alpha_topk_sensitivity}). Third, paired bootstrap comparisons confirm that verification adds a reliable gain (+0.111 in Human Pass@2), while the full-loop wrapper contributes a small, non-decisive increment (+0.021) (Appendix~\ref{app:additional_experiment_tables}, Table~\ref{tab:paired_repair_test}).

Finally, leave-one-perturbation-out and leave-one-source-group-out evaluations confirm that performance drops under held-out operators---particularly missing-event, camera-shake, and temporal-discontinuity perturbations---but remains above prior-only and score-only baselines (Appendix~\ref{app:additional_experiment_tables}, Tables~\ref{tab:leave_one_perturbation}--\ref{tab:leave_one_source_group}). The system does not simply memorize perturbation templates, though unseen-operator generalization is weaker than in-distribution evaluation. Additional checks on adapter executability, real-failure stress slices, and multi-annotator stability are in Appendix~\ref{app:additional_experiment_tables} (Tables~\ref{tab:adapter_executability}--\ref{tab:multi_annotator_stability}).

\section{Discussion}
\label{sec:discussion}

The evidence supports a bounded claim: structured diagnosis records improve failure detection, temporal grounding, and downstream repair over score-only feedback and direct VLM reporting on the current evaluation sets. The controlled fixture isolates known perturbations for auditable comparison; the real-failure subset provides an initial transfer check that confirms the approach degrades under natural artifacts but remains useful.

VideoGenDoctor is best understood as an interface contribution, not a new generator or learned repair algorithm. Its advantages are schema validity, evidence auditability, repair-plan structure, and cost-awareness. Rule+GPT-4V achieves the strongest absolute scores but at higher cost and latency; VideoGenDoctor-full is the recommended default because it preserves the full audit trail at the paper's reference cost point. The modular design allows feature extractors, rules, VLM judges, and patch compilers to be replaced independently.

\paragraph{Limitations.}
The controlled fixture does not represent open-world video diversity. The real-failure subset is failure-enriched---it supports transfer analysis under verified failures but cannot estimate natural prevalence. Automatic pass rates are partly verifier-dependent; we report blind human evaluation as the primary repair metric. Only the \TaxonomyNumCodes{} observed taxonomy codes carry empirical support; the remaining \TaxonomyTotalCodes{} codes are definitional. Repair assumes a downstream generator or editor that consumes at least part of the patch vocabulary. The annotation study uses two adjudicated annotators; broader validation requires more generators, longer videos, additional domains, and larger annotator pools.

\paragraph{Verifier--human agreement.}
Because the closed-loop stopping criterion depends on the system verifier, we measure verifier--human agreement on the 144-video blind-evaluation subset. Agreement between the automatic verifier and the majority human judgment reaches \VerifierHumanAgreement{} (Cohen's $\kappa = \VerifierHumanKappa$), indicating that the automatic pass/fail signal is well-calibrated but not a perfect substitute for human evaluation. The primary Pass@K claims in this paper are therefore based on blind human judgment; automatic Pass@K serves as a scalable proxy for ablation and hyperparameter sweeps.

\paragraph{Break-even analysis.}
At the default operating point, diagnosis costs approximately \$0.24 per video (CLIP + flow + face features, rule engine, and Stage-2 VLM on $K=3$ candidates). Full re-rendering costs approximately \$0.52 per video on current hardware. The break-even failure rate is approximately \BreakEvenFailureRate: when fewer than half of generated videos contain specification violations, diagnosis-guided repair is cheaper than re-rendering everything. When failure rates are higher---as in early-stage prompt engineering or new-domain adaptation---re-render-all may be more economical. This trade-off is visualized in Appendix~\ref{app:moved_main_details}, Figure~\ref{fig:cost_performance_pareto}.

\paragraph{When does it fail?}
Diagnosis degrades most on overlapping artifacts where multiple failure codes are plausible (e.g., camera shake vs. motion jitter), on subtle style shifts below the CLIP embedding sensitivity threshold, and on failures in the last 20\% of the video timeline where fewer keyframes are sampled. The Stage-2 VLM judge occasionally over-corrects when asked about codes outside its effective visual range (e.g., fine-grained shot-type discrimination). False positives on real-normal videos are concentrated among texture-flicker and lighting-shift codes where the rule engine triggers on benign high-frequency variation. These failure modes inform the recommendations for threshold selection and candidate budget in Section~\ref{sec:ablations}.

\section{Conclusion}
VideoGenDoctor replaces scalar video evaluation with auditable diagnosis records that each locate a failure in time, name the affected target, and prescribe a repair action. On a 324-video controlled diagnostic fixture, the default pipeline detects failure codes at macro-F1 0.911, localizes evidence spans at tIoU@0.5 of 0.732, and lifts blind human repair Pass@2 from 27.8\% to 76.4\%. On a 240-video failure-enriched real-generator subset spanning four production generators and 18 taxonomy codes, the pipeline retains macro-F1 0.826 and tIoU@0.5 0.598. The modular design---separating feature extraction, rule-based screening, VLM verification, and patch compilation---allows each component to be replaced independently. These results support structured diagnosis-and-repair as a practical interface for controllable video generation, but do not establish generality beyond the controlled perturbations, generator set, and video lengths evaluated here.

\section{Ethics statement}
\label{sec:ethics}
VideoGenDoctor improves the controllability of generated video through automated diagnosis. Improved generation reliability could be misused to strengthen misleading media. In sensitive deployments, automated repair should be paired with provenance tracking, watermarking, and human review. The failure taxonomy encodes normative judgments about visual quality; creative workflows should expose configurable thresholds and permit human override.

\begin{ack}
This work was supported by the Sichuan Agricultural University College of Information Engineering. We thank the anonymous reviewers for their constructive feedback. The annotation team at Sichuan Agricultural University contributed to the reliability study.

The release package includes the reference implementation, report schema, taxonomy, patch compiler, controlled-fixture generation scripts, annotations, prediction logs, evaluation scripts, VLM baseline prompt templates, and bootstrap confidence-interval scripts. Each run records configuration, model and package versions, random seeds, prompts, raw VLM responses, and output paths. CPU-only recomputation of tables from serialized predictions and annotations completes in seconds; full feature extraction and VLM judging depend on GPU availability and API access.

All newly generated data supporting the findings of this study---including VideoGenDoctor-Bench-v0 (324 controlled perturbation videos, 2,016 human-verified evidence spans), the 240-video failure-enriched real-generator subset (1,536 verified evidence spans), the 120-video real-normal subset, the full 38-code taxonomy, annotation records, prediction logs, and diagnostic reports---will be deposited in a public Zenodo repository upon acceptance and are available for review as part of the supplementary artifact package. The 18 clean source clips used to construct the controlled fixture are obtained from publicly available video datasets under permissive licences. The four production generators (CogVideoX, Wan, Stable Video Diffusion, HunyuanVideo) are publicly available models; checkpoint identifiers are recorded in per-video JSON manifests. No human-subject data or personally identifiable information were collected. The annotation protocol was approved by the institutional review board of Sichuan Agricultural University.
\end{ack}

\FloatBarrier

{\small
\bibliographystyle{plainnat}
\bibliography{references}
}

\appendix


\section{Supplementary Main-Text Results}
\label{app:moved_main_details}
This appendix stores tables and figures moved out of the main text to keep the core narrative compact while preserving the audit trail for comparisons, taxonomy scope, fixture construction, temporal evidence profiles, strengthened VLM baselines, repair visualization, and cost-performance trade-offs.

\begin{table}[t]
\centering
\footnotesize
\caption{Failure types, the evidence signals that detect them, and the repair actions they trigger.}
\label{tab:code_patch_mapping}
\begin{tabular}{p{0.29\textwidth}p{0.33\textwidth}p{0.26\textwidth}}
\toprule
\textbf{Failure type} & \textbf{Evidence signal} & \textbf{Repair action} \\
\midrule
Face identity drift & Face embedding drift across keyframes & Reference-frame injection \\
Body / clothing drift & Character-region appearance drift & Reference-frame or character-anchor injection \\
Camera movement deviation & Camera flow trace or trajectory mismatch & Camera movement override \\
Frozen frames or action stalls & Near-zero flow or stalled pose track & Regenerate affected segment \\
Missing required prop & Object presence check or VLM prop-absence judgment & Prop injection \\
Style inconsistency & Style or color embedding drift across segments & Style correction \\
Background inconsistency & Scene/background embedding drift & Background anchoring \\
\bottomrule
\end{tabular}
\end{table}

\begin{table}[t]
\centering\small
\caption{Dataset statistics for the controlled fixture and failure-enriched real-generator subset. The real-failure split cannot estimate natural failure prevalence.}
\label{tab:dataset_statistics}
\begin{tabular*}{0.95\textwidth}{@{\extracolsep{\fill}}lccccc@{}}
\toprule
\textbf{Split} & \textbf{\#Vid} & \textbf{\#Evidence} & \textbf{AvgDur} & \textbf{PerturbTypes} & \textbf{\#Codes} \\
\midrule
Controlled fixture & 324 & 2016 & 10.46s & 9 & 12 \\
Real-failure subset & 240 & 1536 & 8.93s & -- & 18 \\
\bottomrule
\end{tabular*}
\end{table}

\begin{table}[t]
\centering\small
\caption{Capability comparison with related evaluation and critique paradigms. Here ``Global score'' denotes scalar quality, alignment, or verifier-style pass outputs rather than localized evidence records. Direct VLM reporting is treated as a strong baseline rather than only comparing against score-only benchmarks.}
\label{tab:comparison}
\begin{adjustbox}{max width=\textwidth}
\begin{tabular}{lccccc}
\toprule
\textbf{System} & \textbf{Global score} & \textbf{Code} & \textbf{Evidence} & \textbf{Patch} & \textbf{Closed-Loop} \\
\midrule
VBench & \cmark & \xmark & \xmark & \xmark & \xmark \\
EvalCrafter & \cmark & \xmark & \xmark & \xmark & \xmark \\
VideoScore & \cmark & \xmark & \xmark & \xmark & \xmark \\
T2V-CompBench & \cmark & partial & \xmark & \xmark & \xmark \\
Direct VLM report & partial & partial & partial & partial & \xmark \\
\textbf{VideoGenDoctor} & \cmark & \cmark & \cmark & \cmark & \cmark \\
\bottomrule
\end{tabular}
\end{adjustbox}
\end{table}

\begin{table}[t]
\centering
\small
\caption{Failure-as-code taxonomy groups. Each group defines a namespace for diagnosis records and links to evidence procedures and repair-plan templates.}
\label{tab:failure_taxonomy_groups}
\begin{tabular}{p{0.20\textwidth}p{0.70\textwidth}}
\toprule
\textbf{Group} & \textbf{Operational scope} \\
\midrule
Identity & Identity drift, face/body consistency, wrong character, clothing consistency. \\
Camera & Shot type, camera movement, viewpoint angle, unintended zoom, camera shake. \\
Motion & Action timing, frame continuity, motion physics, frozen frames, segment breaks. \\
Alignment & Prompt adherence, required objects and props, character count, prop placement. \\
Style & Color consistency, texture flicker, compression artifacts, art-style consistency. \\
Scene & Background consistency, lighting, environment match, scene-object stability. \\
\bottomrule
\end{tabular}
\end{table}

\begin{table}[t]
\centering\scriptsize
\caption{Composition of the controlled diagnostic fixture. Each domain group contains three clean source clips. Every source clip is paired with a ShotIR-style specification and perturbed by nine deterministic operators under two random seeds.}
\label{tab:controlled_fixture_composition}
\begin{adjustbox}{max width=\textwidth}
\begin{tabular}{p{0.27\textwidth}cccccc}
\toprule
Domain group & \#Source clips & Avg. duration & Resolution & \#Perturb. ops & \#Seeds & \#Videos \\
\midrule
Indoor / human-object & 3 & 10.1s & 768$\times$432 & 9 & 2 & 54 \\
Outdoor / motion & 3 & 10.8s & 768$\times$432 & 9 & 2 & 54 \\
Multi-object interaction & 3 & 10.4s & 768$\times$432 & 9 & 2 & 54 \\
Camera-control scene & 3 & 11.0s & 768$\times$432 & 9 & 2 & 54 \\
Style / lighting variation & 3 & 10.2s & 768$\times$432 & 9 & 2 & 54 \\
Scene-change / background & 3 & 10.3s & 768$\times$432 & 9 & 2 & 54 \\
\midrule
Total & 18 & 10.46s & -- & 9 & 2 & 324 \\
\bottomrule
\end{tabular}
\end{adjustbox}
\end{table}

\begin{figure}[t]
\centering
\safeincludegraphics[width=\textwidth]{../figures/temporal_evidence_profile.pdf}
\vspace{0.3em}
\safeincludegraphics[width=\textwidth]{../figures/temporal_evidence_real_profile.pdf}
\vspace{0.3em}
\safeincludegraphics[width=\textwidth]{../figures/temporal_evidence_heatmap.pdf}
\caption{Temporal distribution of verified evidence spans in the evaluation sets. Panel (a) aggregates the 2,016 controlled-fixture spans over 2-second bins and shows the front- and mid-loaded pattern expected from deterministic perturbation operators. Panel (b) summarizes the 1,536 real-failure spans over the same time bins, making the later and more diffuse temporal mass directly comparable to panel (a). Panel (c) further breaks the real-failure spans down by failure group and time bin, highlighting that motion and camera failures dominate the later intervals and overlap more strongly than in the controlled fixture.}
\label{fig:temporal_evidence_distribution}
\end{figure}

\begin{table}[t]
\centering\scriptsize
\caption{Strengthened VLM-agent baselines on the controlled fixture. All rows are evaluated under the same downstream repair protocol; diagnosis-only front ends are explicitly paired with that shared repair loop when human repair success is reported. Self-consistency improves direct VLM reporting but increases cost. The validity column reports whether the emitted action stays inside the supported repair-plan vocabulary; it does not claim generator-adapter executability. Best values are boldfaced; higher is better except for relative cost, where lower is better.}
\label{tab:strong_vlm_agent_baselines}
\begin{adjustbox}{max width=\textwidth}
\begin{tabular}{lcccccc}
\toprule
Method & Macro-F1 & tIoU@0.5 & Human Pass@2 & Schema validity & Patch vocabulary validity & Relative cost \\
\midrule
VLM free-form report & 0.8402 & 0.5986 & 0.6125 & 0.742 & 0.604 & 1.12$\times$ \\
VLM structured report & 0.8869 & 0.6639 & 0.6810 & 0.914 & 0.732 & 1.16$\times$ \\
VLM structured + self-consistency & 0.8978 & 0.6824 & 0.7042 & 0.951 & 0.755 & 3.48$\times$ \\
VLM temporal proposal + patch & 0.8915 & 0.7018 & 0.7076 & 0.936 & 0.802 & 2.05$\times$ \\
VideoGenDoctor-full & 0.9110 & 0.7316 & 0.7639 & 0.991 & 0.953 & \textbf{1.00}$\times$ \\
Rule+GPT-4V + repair loop & \textbf{0.9276} & \textbf{0.7688} & \textbf{0.7917} & \textbf{0.993} & \textbf{0.957} & 2.75$\times$ \\
\bottomrule
\end{tabular}
\end{adjustbox}
\end{table}

\begin{figure}[t]
\centering
\safeincludegraphics[width=\textwidth]{../figures/manuscript_fig3.pdf}
\caption{Closed-loop repair comparison on the controlled fixture. Left: automatic verifier Pass@1/2 on all 324 controlled videos. Right: blind human Pass@1/2 on the stratified 144-video human-evaluation subset. Automatic pass rates may be coupled to the system verifier, while human pass rates provide an independent check of repair success.}
\label{fig:repair_curve}
\end{figure}

\begin{table}[t]
\centering\scriptsize
\caption{Cost-performance trade-off under a shared downstream repair protocol. Diagnosis-only front ends are paired with the same repair loop before reporting Human Pass@2. Cost is reported relative to the default VideoGenDoctor-full configuration. Values in the last two columns are multiplicative factors relative to VideoGenDoctor-full (1.00$\times$). Best values are boldfaced; higher is better for task metrics, lower is better for cost and latency.}
\label{tab:cost_performance}
\begin{adjustbox}{max width=\textwidth}
\begin{tabular}{p{0.28\textwidth}ccccc}
\toprule
Method & Macro-F1 & tIoU@0.5 & Human Pass@2 & Relative cost & Relative latency \\
\midrule
Rule-only + repair loop & 0.8926 & 0.6714 & 0.5988 & \textbf{0.42}$\times$ & \textbf{0.58}$\times$ \\
Rule+Open-VLM + repair loop & 0.9047 & 0.7062 & 0.7160 & 0.73$\times$ & 0.87$\times$ \\
VideoGenDoctor-full & 0.9110 & 0.7316 & 0.7639 & 1.00$\times$ & 1.00$\times$ \\
Rule+GPT-4V + repair loop & \textbf{0.9276} & \textbf{0.7688} & \textbf{0.7917} & 2.75$\times$ & 2.31$\times$ \\
VLM structured report & 0.8869 & 0.6639 & 0.6810 & 1.16$\times$ & 1.20$\times$ \\
\bottomrule
\end{tabular}
\end{adjustbox}
\end{table}

\begin{figure}[t]
\centering
\safeincludegraphics[width=\textwidth]{../figures/cost_performance_pareto.pdf}
\caption{Pareto view of the cost-performance trade-off under the shared downstream repair protocol. The default VideoGenDoctor-full operating point improves human repair success over lower-cost rule variants and direct structured VLM reporting, while the higher-cost Rule+GPT-4V reference obtains the strongest Human Pass@2 at substantially higher relative cost.}
\label{fig:cost_performance_pareto}
\end{figure}


% Begin inlined from appendix.tex
\section{Additional Experimental Tables}
\label{app:additional_experiment_tables}
This appendix collects construction details, audit tables, stress analyses, per-slice results, and visualizations that support the main experimental claims without interrupting the main narrative.

\subsection{Diagnosis and Repair Visualizations}
\label{app:diagnosis_repair_viz}

\begin{figure}[t]
\centering
\safeincludegraphics[width=\textwidth]{../figures/controlled_vs_real_paired.pdf}
\caption{Controlled-to-real transfer across four metrics. Solid bars: controlled fixture; hatched bars: real-failure subset. The degradation gap is largest for tIoU@0.5 and smallest for patch usefulness.}
\label{fig:controlled_vs_real_paired}
\end{figure}

\begin{figure}[t]
\centering
\safeincludegraphics[width=\textwidth]{../figures/feature_ablation_horizontal.pdf}
\caption{Per-feature ablation horizontal bar chart. Optical flow is the single most critical signal; single-signal configurations all underperform substantially.}
\label{fig:feature_ablation}
\end{figure}

\begin{figure}[t]
\centering
\safeincludegraphics[width=\textwidth]{../figures/repair_ablation_bars.pdf}
\caption{Repair ablation as grouped bars. Human Pass@2 and patch usefulness rise additively as code, span, evidence, and verification are incorporated.}
\label{fig:repair_ablation_bars}
\end{figure}

\begin{table}[t]
\centering\scriptsize
\caption{Threshold sensitivity. $\tau=0.50$ balances diagnosis coverage against false positives and unnecessary patching.}
\label{tab:threshold_sensitivity}
\begin{tabular}{ccccc}
\toprule
Threshold $\tau$ & Macro-F1 & tIoU@0.5 & Real-normal FPR & Unnecessary patch rate \\
\midrule
0.30 & 0.902 & 0.716 & 0.158 & 0.142 \\
0.40 & 0.910 & 0.728 & 0.125 & 0.111 \\
0.50 & \textbf{0.911} & \textbf{0.732} & 0.100 & 0.087 \\
0.60 & 0.904 & 0.719 & 0.083 & 0.071 \\
0.70 & 0.886 & 0.694 & \textbf{0.067} & \textbf{0.058} \\
\bottomrule
\end{tabular}
\end{table}

\begin{figure}[t]
\centering
\safeincludegraphics[width=\textwidth]{../figures/threshold_sensitivity.pdf}
\caption{Threshold sensitivity curves. Dashed line at $\tau=0.50$ marks the operating point where macro-F1 and tIoU@0.5 peak.}
\label{fig:threshold_sensitivity}
\end{figure}

\begin{figure}[t]
\centering
\safeincludegraphics[width=\textwidth]{../figures/alpha_topk_sensitivity.pdf}
\caption{Sensitivity to judge weight $\alpha$ and candidate budget $K$. $\alpha=0.6$, $K=3$ balances diagnosis accuracy and repair success.}
\label{fig:alpha_topk_sensitivity}
\end{figure}

\begin{table}[t]
\centering\scriptsize
\caption{Paired bootstrap differences in blind human Pass@2. Verification adds a reliable gain; the full-loop wrapper adds a small, non-decisive increment.}
\label{tab:paired_repair_test}
\begin{tabular}{lccc}
\toprule
Comparison & $\Delta$ Human Pass@2 & 95\% CI & Interpretation \\
\midrule
Patch-only vs Score-only & 0.3541 & [0.250, 0.458] & substantial gain \\
VLM-to-patch vs Patch-only & 0.0417 & [-0.054, 0.139] & small / not decisive \\
Patch+Judge vs Patch-only & 0.1112 & [0.021, 0.215] & verification helps \\
VideoGenDoctor-full vs Patch+Judge & 0.0208 & [-0.063, 0.104] & small / not decisive \\
\bottomrule
\end{tabular}
\end{table}

\subsection{Extended Experiment Analyses}
\label{app:experiment_extensions}

\paragraph{Adapter executability.}
Table~\ref{tab:adapter_executability} and Figure~\ref{fig:adapter_executability} separate repair-plan quality from backend executability by grouping generated plans into L0--L3 adapter levels. The trend is consistent with the intended repair interface: L2/L3 actions are more often executable and yield higher Human Pass@2, while L0/L1 outputs remain useful as conservative abstentions or prompt-level repair instructions but provide weaker direct repair success.

\begin{table}[t]
\centering\scriptsize
\caption{Adapter-executability stratification. L0: abstention; L1: prompt-level instructions; L2: partially executable actions; L3: native generator/editor controls.}
\label{tab:adapter_executability}
\begin{tabular}{lcccc}
\toprule
Adapter level & Plan share & Adapter-executable rate & Human Pass@2 & New artifacts \\
\midrule
L0 abstain & 0.100 & 0.000 & 0.280 & 0.040 \\
L1 plan & 0.310 & 0.280 & 0.540 & 0.100 \\
L2 partial & 0.420 & 0.740 & 0.730 & 0.140 \\
L3 native & 0.170 & 0.960 & 0.810 & 0.160 \\
\bottomrule
\end{tabular}
\end{table}

\begin{figure}[t]
\centering
\safeincludegraphics[width=\textwidth]{../figures/adapter_executability.pdf}
\caption{Adapter-executability stratification by repair-plan level. The left panel reports the share of generated plans by adapter level; the right panel compares adapter-executable rate, Human Pass@2, and new-artifact rate within each level.}
\label{fig:adapter_executability}
\end{figure}

\paragraph{Stronger real-failure stress slices.}
Table~\ref{tab:extended_real_validation} and Figure~\ref{fig:extended_real_validation} extend the real-failure check to harder slices: additional generators, longer videos, style-shifted prompts, and more complex prompt specifications. The stress slices degrade relative to the base real-failure subset, especially for long videos and complex prompts, but the retained performance remains above the score-only and prior-only operating ranges reported in the main experiments.

\begin{table}[t]
\centering\scriptsize
\caption{Stress validation on harder real-failure slices. Intervals are bootstrap half-widths for Figure~\ref{fig:extended_real_validation}.}
\label{tab:extended_real_validation}
\begin{tabular}{lcccc}
\toprule
Slice & Macro-F1 & tIoU@0.5 & Human Pass@2 & FPR \\
\midrule
Base real-failure & 0.826 $\pm$ 0.020 & 0.598 $\pm$ 0.017 & 0.691 $\pm$ 0.028 & 0.100 $\pm$ 0.010 \\
New generators & 0.802 $\pm$ 0.024 & 0.574 $\pm$ 0.020 & 0.662 $\pm$ 0.034 & 0.126 $\pm$ 0.012 \\
Long videos & 0.784 $\pm$ 0.029 & 0.552 $\pm$ 0.025 & 0.628 $\pm$ 0.041 & 0.142 $\pm$ 0.014 \\
Style shift & 0.811 $\pm$ 0.025 & 0.581 $\pm$ 0.022 & 0.674 $\pm$ 0.036 & 0.118 $\pm$ 0.013 \\
Complex prompts & 0.793 $\pm$ 0.027 & 0.560 $\pm$ 0.023 & 0.641 $\pm$ 0.039 & 0.135 $\pm$ 0.014 \\
\bottomrule
\end{tabular}
\end{table}

\begin{figure}[t]
\centering
\safeincludegraphics[width=\textwidth]{../figures/extended_real_validation.pdf}
\caption{Stress validation on harder real-failure slices. Left: diagnosis, localization, and human repair with uncertainty. Right: false-positive rate on matched real-normal controls.}
\label{fig:extended_real_validation}
\end{figure}

\paragraph{Multi-annotator stability.}
Table~\ref{tab:multi_annotator_stability} and Figure~\ref{fig:multi_annotator_stability} extend the annotation reliability analysis from dual annotation to a three-annotator setting. Code-level agreement remains substantial, while span agreement is lower on real-failure samples, consistent with the higher ambiguity of naturally occurring failures and less isolated temporal evidence.

\begin{table}[t]
\centering\scriptsize
\caption{Three-annotator stability analysis. Boundary disagreement in seconds (lower is better).}
\label{tab:multi_annotator_stability}
\begin{tabular}{lccccc}
\toprule
Subset & \#Videos & \#Candidate spans & Annotators & Fleiss' $\kappa$ & Mean pairwise tIoU / Boundary disagreement \\
\midrule
Controlled & 72 & 468 & 3 & 0.842 & 0.704 / 0.38s \\
Real-failure & 48 & 300 & 3 & 0.782 & 0.628 / 0.52s \\
Combined & 120 & 768 & 3 & 0.813 & 0.666 / 0.45s \\
\bottomrule
\end{tabular}
\end{table}

\begin{figure}[t]
\centering
\safeincludegraphics[width=\textwidth]{../figures/multi_annotator_stability.pdf}
\caption{Multi-annotator stability. Code-level $\kappa$ exceeds temporal-boundary agreement; the real-failure split shows highest ambiguity.}
\label{fig:multi_annotator_stability}
\end{figure}

\begin{table}[t]
\centering\scriptsize
\caption{Operator-template alignment analysis for the controlled fixture. The table makes explicit which observable cues, rule signals, failure codes, and repair-plan templates are aligned by construction.}
\label{tab:operator_alignment}
\begin{adjustbox}{max width=\textwidth}
\begin{tabular}{p{0.18\textwidth}p{0.16\textwidth}p{0.16\textwidth}p{0.16\textwidth}p{0.17\textwidth}p{0.11\textwidth}}
\toprule
Perturbation operator & Observable cue & Rule / feature signal & Failure code & Repair-plan template & Alignment type \\
\midrule
Identity-anchor removal & face or body appearance drift & face embedding drift; character-region embedding drift & \tblcode{ID\_FACE\_DRIFT}, \tblcode{ID\_BODY\_DRIFT} & \tblcode{inject\_reference\_frame}; \tblcode{replace\_character\_anchor} & identity-anchor \\
Required-prop dropping & required object becomes absent or inconsistent & object detector absence; VLM prop-absence judgment & \tblcode{AL\_PROP\_MISSING} & \tblcode{inject\_prop} & prop-presence \\
Camera-movement alteration & requested trajectory is replaced by static or wrong motion & global flow direction; trajectory mismatch & \tblcode{CA\_MOVE\_WRONG} & \tblcode{set\_camera\_movement} & motion-control \\
Shot-type / framing alteration & framing violates requested shot type & crop statistics; VLM framing judgment & \tblcode{CA\_SHOT\_TYPE\_WRONG} & \tblcode{adjust\_framing} & framing-control \\
Camera-shake injection & unstable camera with unintended jitter & flow instability; camera-shake heuristic & \tblcode{CA\_SHAKE} & \tblcode{stabilize\_camera} & stability-control \\
Action dropping / shifting & required event missing or misordered & action-track gap; VLM event-order judgment & \tblcode{MO\_EVENT\_MISSING} & \tblcode{inject\_action\_keyframe}; \tblcode{regenerate\_segment} & event-structure \\
Temporal jitter / frame-drop injection & discontinuity, duplicate frames, or rate mismatch & flow anomaly; near-zero motion plateau & \tblcode{MO\_JITTER}, \tblcode{MO\_SEGMENT\_BREAK} & \tblcode{normalize\_frame\_rate}; \tblcode{regenerate\_segment} & temporal-continuity \\
Compression re-encoding & blockiness and deblocking artifacts & compression detector; texture degradation signal & \tblcode{ST\_COMPRESSION\_ARTIFACT} & \tblcode{apply\_enhancement} & quality-restoration \\
Background-anchor weakening & scene background drifts across the shot & scene embedding drift; background-anchor mismatch & \tblcode{SC\_BG\_INCONSISTENCY} & \tblcode{fix\_background\_anchor} & scene-anchor \\
\bottomrule
\end{tabular}
\end{adjustbox}
\end{table}

\begin{table}[t]
\centering\scriptsize
\caption{Source composition of the failure-enriched real-generator subset. We generate 480 raw videos and retain 240 samples with annotator-verified visible failures. Sampling is balanced at 60 retained failures per generator, and generator-specific prompt or conditioning adapters are logged in the construction manifest. The subset is used only as an initial transfer check, not to estimate natural failure prevalence. We use publicly identifiable generator checkpoints to make the subset construction auditable and reproducible.}
\label{tab:real_failure_source}
\begin{adjustbox}{max width=\textwidth}
\begin{tabular}{p{0.16\textwidth}p{0.18\textwidth}p{0.22\textwidth}cccc}
\toprule
Generator & Version / checkpoint & Input type & \#Prompts / specs & \#Videos & Avg. duration & \#Unique observed codes \\
\midrule
CogVideoX & CogVideoX1.5-5B & text / ShotIR-to-prompt & 60 & 60 & 8.4s & 15 \\
Wan & Wan2.1-T2V-14B & text / ShotIR-to-prompt & 60 & 60 & 9.2s & 16 \\
Stable Video Diffusion & SVD-XT-1.1 & image-to-video / reference frame & 60 & 60 & 8.6s & 14 \\
HunyuanVideo & HunyuanVideo-1.5 & text / multi-shot prompt & 60 & 60 & 9.5s & 17 \\
\midrule
Total & -- & -- & 240 & 240 & 8.93s & 18 \\
\bottomrule
\end{tabular}
\end{adjustbox}
\end{table}

\begin{table}[t]
\centering\scriptsize
\caption{Real-failure and real-normal evaluation. The real-normal subset contains 120 videos retained from the same 480 raw generations but judged to contain no obvious failure. Whenever a diagnosis-oriented front end is listed, patch-related columns are computed after pairing that front end with the shared downstream repair loop. Best values are boldfaced; higher is better except for real-normal FPR and unnecessary patch rate, where lower is better.}
\label{tab:real_normal_control}
\begin{adjustbox}{max width=\textwidth}
\begin{tabular}{lccccc}
\toprule
Method & Real-failure Recall & Real-normal FPR & Unnecessary Patch Rate & No-op Accuracy & Specificity \\
\midrule
Score-only & 0.621 & 0.258 & 0.231 & 0.769 & 0.742 \\
Rule-only & 0.806 & 0.168 & 0.154 & 0.846 & 0.832 \\
VLM structured report & 0.818 & 0.217 & 0.205 & 0.795 & 0.783 \\
VLM structured + self-consistency & 0.837 & 0.167 & 0.150 & 0.850 & 0.833 \\
VLM temporal proposal + patch & 0.848 & 0.192 & 0.183 & 0.817 & 0.808 \\
VideoGenDoctor-diagnosis & \textbf{0.862} & 0.108 & 0.092 & 0.908 & 0.892 \\
VideoGenDoctor-full & 0.858 & \textbf{0.100} & \textbf{0.087} & \textbf{0.913} & \textbf{0.900} \\
\bottomrule
\end{tabular}
\end{adjustbox}
\end{table}

\begin{table}[t]
\centering\scriptsize
\caption{Stress analysis on 72 ambiguous specification-violation cases. These samples contain possible but non-adjudicable deviations from the specification. Whenever a diagnosis-oriented front end is listed, patch-related columns are computed after pairing that front end with the shared downstream repair loop. Best values are boldfaced; higher is better for abstention-oriented columns, lower is better for false concrete patch, wrong-target patch, over-repair, and new artifacts.}
\label{tab:ambiguous_violation_cases}
\begin{adjustbox}{max width=\textwidth}
\begin{tabular}{lcccccc}
\toprule
Method & Abstention rate & Appropriate abstention & False concrete patch & Wrong-target patch & Over-repair & New artifacts \\
\midrule
Score-only & 0.292 & 0.361 & 0.514 & 0.222 & 0.347 & 0.083 \\
VLM structured report & 0.181 & 0.292 & 0.611 & 0.278 & 0.431 & 0.153 \\
VLM temporal proposal + patch & 0.222 & 0.347 & 0.556 & 0.236 & 0.389 & 0.194 \\
VideoGenDoctor-diagnosis & \textbf{0.500} & \textbf{0.639} & 0.292 & 0.125 & 0.208 & \textbf{0.069} \\
VideoGenDoctor-full & 0.472 & 0.625 & \textbf{0.250} & \textbf{0.111} & \textbf{0.181} & 0.111 \\
\bottomrule
\end{tabular}
\end{adjustbox}
\end{table}

\begin{table}[t]
\centering\scriptsize
\caption{Stress analysis on 48 low-quality or non-adjudicable cases. These samples are dominated by global degradation, unclear content, or insufficient evidence for reliable code-level labels. Whenever a diagnosis-oriented front end is listed, patch-related columns are computed after pairing that front end with the shared downstream repair loop. Best values are boldfaced; higher is better for quality-gated abstention, lower is better for the remaining columns.}
\label{tab:low_quality_cases}
\begin{adjustbox}{max width=\textwidth}
\begin{tabular}{lccccc}
\toprule
Method & Quality-gated abstention & False concrete patch & Wrong-target patch & Over-repair & New artifacts \\
\midrule
Score-only & 0.417 & 0.375 & 0.125 & 0.208 & 0.063 \\
VLM structured report & 0.292 & 0.458 & 0.188 & 0.271 & 0.104 \\
VLM temporal proposal + patch & 0.333 & 0.396 & 0.167 & 0.250 & 0.146 \\
VideoGenDoctor-diagnosis & \textbf{0.667} & 0.146 & 0.063 & 0.104 & \textbf{0.042} \\
VideoGenDoctor-full & 0.625 & \textbf{0.125} & \textbf{0.042} & \textbf{0.083} & 0.083 \\
\bottomrule
\end{tabular}
\end{adjustbox}
\end{table}

\begin{table}[t]
\centering\scriptsize
\caption{Observed failure-code distribution in the controlled fixture and real-failure subset. The full taxonomy contains 38 codes, but empirical claims are restricted to the observed codes below.}
\label{tab:code_distribution}
\begin{adjustbox}{max width=\textwidth}
\begin{tabular}{p{0.30\textwidth}p{0.14\textwidth}ccc}
\toprule
Failure code & Group & Controlled spans & Real-failure spans & Total spans \\
\midrule
\tblcode{ID\_FACE\_DRIFT} & Identity & 150 & 112 & 262 \\
\tblcode{ID\_BODY\_DRIFT} & Identity & 132 & 76 & 208 \\
\tblcode{ID\_CLOTHING\_CHANGE} & Identity & -- & 60 & 60 \\
\tblcode{CA\_MOVE\_WRONG} & Camera & 240 & 104 & 344 \\
\tblcode{CA\_SHOT\_TYPE\_WRONG} & Camera & 144 & 84 & 228 \\
\tblcode{CA\_SHAKE} & Camera & 126 & 70 & 196 \\
\tblcode{CA\_WRONG\_ANGLE} & Camera & -- & 68 & 68 \\
\tblcode{MO\_JITTER} & Motion & 174 & 124 & 298 \\
\tblcode{MO\_FROZEN\_FRAME} & Motion & 156 & 72 & 228 \\
\tblcode{MO\_SEGMENT\_BREAK} & Motion & 162 & 88 & 250 \\
\tblcode{MO\_EVENT\_MISSING} & Motion & 126 & 90 & 216 \\
\tblcode{MO\_ACTION\_ORDER\_WRONG} & Motion & -- & 56 & 56 \\
\tblcode{AL\_PROP\_MISSING} & Alignment & 198 & 116 & 314 \\
\tblcode{AL\_TEXT\_PROMPT\_MISMATCH} & Alignment & -- & 132 & 132 \\
\tblcode{ST\_COMPRESSION\_ARTIFACT} & Style & 252 & 80 & 332 \\
\tblcode{ST\_COLOR\_SHIFT} & Style & -- & 56 & 56 \\
\tblcode{SC\_BG\_INCONSISTENCY} & Scene & 156 & 86 & 242 \\
\tblcode{SC\_OBJECT\_TELEPORT} & Scene & -- & 62 & 62 \\
\midrule
Total & -- & 2016 & 1536 & 3552 \\
\bottomrule
\end{tabular}
\end{adjustbox}
\end{table}

\begin{table}[t]
\centering\scriptsize
\caption{Annotation reliability. Dual annotation is performed on a reliability subset before adjudication.}
\label{tab:annotation_reliability}
\begin{adjustbox}{max width=\textwidth}
\begin{tabular}{p{0.28\textwidth}ccccc}
\toprule
Subset & \#Videos & \#Candidate spans & \#Annotators & Code agreement & Span agreement \\
\midrule
Controlled reliability subset & 72 & 468 & 2 & $\kappa$ = 0.858 & tIoU = 0.716 \\
Real-failure reliability subset & 48 & 300 & 2 & $\kappa$ = 0.803 & tIoU = 0.641 \\
Combined & 120 & 768 & 2 & $\kappa$ = 0.831 & tIoU = 0.681 \\
\bottomrule
\end{tabular}
\end{adjustbox}
\end{table}

\begin{table}[t]
\centering\scriptsize
\caption{Implementation details for VLM-based judges and external diagnosis/repair baselines.}
\label{tab:vlm_details}
\begin{adjustbox}{max width=\textwidth}
\begin{tabular}{>{\raggedright\arraybackslash}p{0.24\textwidth}>{\raggedright\arraybackslash}p{0.19\textwidth}cc>{\raggedright\arraybackslash}p{0.15\textwidth}>{\raggedright\arraybackslash}p{0.14\textwidth}}
\toprule
Setting & Model & Frames & Res. & Prompt & Rel. cost \\
\midrule
Rule+Open-VLM & Qwen2-VL-7B-Instruct & 18 & 512px short side & structured QA & 0.34$\times$ \\
Rule+GPT-4V & GPT-4V endpoint & 18 & 512px short side & structured QA & 2.75$\times$ \\
VLM free-form report & GPT-4V endpoint & 12 & 512px short side & free-form critique & 1.12$\times$ \\
VLM structured report & GPT-4V endpoint & 12 & 512px short side & JSON schema & 1.16$\times$ \\
VLM structured + self-consistency & GPT-4V endpoint & 12 frames $\times$ 3 runs & 512px short side & JSON schema, 3 runs & 3.48$\times$ \\
VLM temporal proposal + patch & GPT-4V endpoint & 12 & 512px short side & span proposal + JSON & 2.05$\times$ \\
VLM-to-patch & GPT-4V endpoint & 12 & 512px short side & JSON patch & 1.19$\times$ \\
Score+LLM-to-patch & scalar evaluator + LLM & 0 & -- & score-to-patch JSON & N/A \\
\bottomrule
\end{tabular}
\end{adjustbox}
\vspace{2pt}
\begin{minipage}{0.98\linewidth}
\footnotesize\emph{Note.} Most direct VLM baselines share the same hosted GPT-4V endpoint~\cite{gpt4v2023} to isolate prompt- and schema-level effects rather than backbone changes; Rule+Open-VLM provides an open-model reference. Hosted endpoint calls are logged with raw responses for auditability. `N/A' denotes a scalar-evaluator-dominated pipeline whose cost is not directly comparable to frame-based VLM calls.
\end{minipage}
\end{table}

\begin{table}[t]
\centering\scriptsize
\caption{Per-code reliability on the real-failure subset. The subset is more challenging than the controlled fixture, especially for motion, camera, and scene-level failures.}
\label{tab:real_per_code}
\begin{adjustbox}{max width=\textwidth}
\begin{tabular}{p{0.31\textwidth}ccccc}
\toprule
Failure code & \#Spans & Precision & Recall & F1 & tIoU@0.5 \\
\midrule
\tblcode{ID\_FACE\_DRIFT} & 112 & 0.884 & 0.848 & 0.866 & 0.648 \\
\tblcode{ID\_BODY\_DRIFT} & 76 & 0.856 & 0.823 & 0.839 & 0.615 \\
\tblcode{ID\_CLOTHING\_CHANGE} & 60 & 0.834 & 0.798 & 0.816 & 0.586 \\
\tblcode{CA\_MOVE\_WRONG} & 104 & 0.866 & 0.832 & 0.849 & 0.628 \\
\tblcode{CA\_SHOT\_TYPE\_WRONG} & 84 & 0.838 & 0.803 & 0.820 & 0.590 \\
\tblcode{CA\_SHAKE} & 70 & 0.820 & 0.779 & 0.799 & 0.559 \\
\tblcode{CA\_WRONG\_ANGLE} & 68 & 0.825 & 0.786 & 0.805 & 0.568 \\
\tblcode{MO\_JITTER} & 124 & 0.823 & 0.777 & 0.799 & 0.565 \\
\tblcode{MO\_FROZEN\_FRAME} & 72 & 0.899 & 0.854 & 0.876 & 0.663 \\
\tblcode{MO\_SEGMENT\_BREAK} & 88 & 0.827 & 0.792 & 0.809 & 0.582 \\
\tblcode{MO\_EVENT\_MISSING} & 90 & 0.807 & 0.762 & 0.784 & 0.542 \\
\tblcode{MO\_ACTION\_ORDER\_WRONG} & 56 & 0.792 & 0.748 & 0.769 & 0.524 \\
\tblcode{AL\_PROP\_MISSING} & 116 & 0.893 & 0.863 & 0.878 & 0.669 \\
\tblcode{AL\_TEXT\_PROMPT\_MISMATCH} & 132 & 0.842 & 0.794 & 0.817 & 0.596 \\
\tblcode{ST\_COMPRESSION\_ARTIFACT} & 80 & 0.914 & 0.874 & 0.893 & 0.688 \\
\tblcode{ST\_COLOR\_SHIFT} & 56 & 0.846 & 0.806 & 0.825 & 0.601 \\
\tblcode{SC\_BG\_INCONSISTENCY} & 86 & 0.832 & 0.795 & 0.813 & 0.578 \\
\tblcode{SC\_OBJECT\_TELEPORT} & 62 & 0.823 & 0.782 & 0.802 & 0.545 \\
\bottomrule
\end{tabular}
\end{adjustbox}
\end{table}

\begin{table}[t]
\centering\scriptsize
\caption{Per-generator results on the real-failure and real-normal subsets. FPR is computed on 30 real-normal samples per generator.}
\label{tab:per_generator_results}
\begin{adjustbox}{max width=\textwidth}
\begin{tabular}{lcccclc}
\toprule
Generator & Macro-F1 & tIoU@0.5 & Human Pass@2 & Real-normal FPR & Main failure groups & Normal samples \\
\midrule
CogVideoX & 0.834 & 0.609 & 0.702 & 0.092 & Camera, Motion, Alignment & 30 \\
Wan & 0.821 & 0.588 & 0.681 & 0.108 & Motion, Camera, Identity & 30 \\
Stable Video Diffusion & 0.842 & 0.617 & 0.696 & 0.083 & Identity, Style, Scene & 30 \\
HunyuanVideo & 0.809 & 0.579 & 0.684 & 0.117 & Camera, Motion, Alignment & 30 \\
\midrule
Average & 0.827 & 0.598 & 0.691 & 0.100 & -- & 120 \\
\bottomrule
\end{tabular}
\end{adjustbox}
\end{table}

\begin{table}[t]
\centering\scriptsize
\caption{Leave-one-perturbation-out evaluation. Each held-out operator is excluded from threshold tuning and then used only for testing.}
\label{tab:leave_one_perturbation}
\begin{adjustbox}{max width=\textwidth}
\begin{tabular}{lccc}
\toprule
Held-out perturbation operator & Macro-F1 & tIoU@0.5 & Top-3 hit \\
\midrule
Identity-anchor removal & 0.884 & 0.681 & 0.644 \\
Required-prop dropping & 0.902 & 0.704 & 0.672 \\
Camera-movement alteration & 0.871 & 0.653 & 0.612 \\
Shot-type / framing alteration & 0.862 & 0.641 & 0.606 \\
Camera-shake injection & 0.851 & 0.628 & 0.590 \\
Action dropping / shifting & 0.838 & 0.612 & 0.568 \\
Temporal jitter / frame-drop injection & 0.846 & 0.621 & 0.577 \\
Compression re-encoding & 0.910 & 0.719 & 0.688 \\
Background-anchor weakening & 0.868 & 0.646 & 0.610 \\
\midrule
Average & 0.870 & 0.656 & 0.619 \\
\bottomrule
\end{tabular}
\end{adjustbox}
\end{table}

\begin{table}[t]
\centering\scriptsize
\caption{Leave-one-source-group-out evaluation over the six controlled-fixture domain groups.}
\label{tab:leave_one_source_group}
\begin{tabular}{lccc}
\toprule
Held-out source group & Macro-F1 & tIoU@0.5 & Top-3 hit \\
\midrule
Indoor / human-object & 0.899 & 0.714 & 0.665 \\
Outdoor / motion & 0.887 & 0.702 & 0.651 \\
Multi-object interaction & 0.892 & 0.706 & 0.657 \\
Camera-control scene & 0.875 & 0.683 & 0.631 \\
Style / lighting variation & 0.907 & 0.722 & 0.676 \\
Scene-change / background & 0.884 & 0.694 & 0.642 \\
\midrule
Average & 0.891 & 0.704 & 0.654 \\
\bottomrule
\end{tabular}
\end{table}

\section{Full Failure Taxonomy}
\label{app:taxonomy}
Table~\ref{tab:full_taxonomy} reports the machine-readable taxonomy used by
VideoGenDoctor v0.1. The v0.1 taxonomy contains \TaxonomyTotalCodes{} codes
grouped into six categories. Each code is paired with an operational definition,
an evidence procedure, and a default patch template. The current benchmark
evaluates the \TaxonomyNumCodes{} codes observed in VideoGenDoctor-Bench-v0,
while the remaining codes define the extensible namespace used by the report
schema and patch compiler.

\begingroup
\scriptsize
\setlength{\tabcolsep}{2pt}
\renewcommand{\arraystretch}{1.08}
\begin{longtable}{
  >{\raggedright\arraybackslash}p{0.23\textwidth}
  >{\raggedright\arraybackslash}p{0.27\textwidth}
  >{\raggedright\arraybackslash}p{0.21\textwidth}
  >{\raggedright\arraybackslash}p{0.13\textwidth}}
\caption{Full VideoGenDoctor v0.1 failure taxonomy used by the report schema and patch compiler.}
\label{tab:full_taxonomy}\\
\toprule
\textbf{Code} & \textbf{Definition} & \textbf{Evidence Procedure} & \textbf{Patch Template} \\
\midrule
\endfirsthead
\toprule
\textbf{Code} & \textbf{Definition} & \textbf{Evidence Procedure} & \textbf{Patch Template} \\
\midrule
\endhead
\midrule
\multicolumn{4}{r}{Continued on next page} \\
\endfoot
\bottomrule
\endlastfoot
\multicolumn{4}{l}{Identity} \\
\tblcode{ID\_FACE\_DRIFT} & Face identity changes inconsistently between frames or segments. & Compare face embeddings across keyframes and flag large cosine distance. & \tblcode{inject\_reference\_frame} \\
\tblcode{ID\_BODY\_DRIFT} & Body appearance, clothing, or build changes inconsistently. & Compare character-region embeddings across segments. & \tblcode{inject\_reference\_frame} \\
\tblcode{ID\_WRONG\_CHARACTER} & A visible character does not match any specified character. & Check face embeddings against known character anchors. & \tblcode{replace\_character\_anchor} \\
\tblcode{ID\_CLOTHING\_CHANGE} & Clothing changes between shots without narrative justification. & Measure character-region appearance drift across segments. & \tblcode{fix\_character\_consistency} \\
\tblcode{ID\_MULTIPLE\_FACES} & Multiple distinct faces appear when one character is expected. & Count distinct identities detected in a single-character shot. & \tblcode{adjust\_character\_count} \\
\tblcode{ID\_NO\_FACE\_VISIBLE} & An expected face is absent when the shot requires face visibility. & Check face detection under close-up or medium-shot constraints. & \tblcode{adjust\_framing} \\
\midrule
\multicolumn{4}{l}{Scene} \\
\tblcode{SC\_BG\_INCONSISTENCY} & The background changes unexpectedly between segments. & Measure background-region embedding drift across consecutive segments. & \tblcode{fix\_background\_anchor} \\
\tblcode{SC\_ENVIRONMENT\_MISMATCH} & The scene environment does not match the specified location. & Compare frames with the location description using image-text similarity. & \tblcode{set\_environment} \\
\tblcode{SC\_LIGHTING\_SHIFT} & Lighting changes inconsistently within a shot. & Measure per-frame luminance variation within the segment. & \tblcode{normalize\_lighting} \\
\tblcode{SC\_OBJECT\_TELEPORT} & A scene object appears or disappears without plausible motion. & Track object boxes across adjacent frames. & \tblcode{stabilize\_scene\_objects} \\
\tblcode{SC\_BG\_FLICKER} & The background flickers or pulses unnaturally. & Detect high-frequency luminance variation in background regions. & \tblcode{fix\_background\_anchor} \\
\tblcode{SC\_WRONG\_TIME\_OF\_DAY} & The time of day does not match the specification. & Classify sky and ambient lighting against the expected time. & \tblcode{set\_time\_of\_day} \\
\midrule
\multicolumn{4}{l}{Motion} \\
\tblcode{MO\_JITTER} & Unnatural frame-to-frame jitter appears outside intended camera motion. & Measure optical-flow magnitude variance and direction instability. & \tblcode{normalize\_frame\_rate} \\
\tblcode{MO\_FRAME\_DROP} & Duplicate or dropped frames produce visible stutter. & Detect near-zero adjacent-frame flow when motion is expected. & \tblcode{interpolate\_frames} \\
\tblcode{MO\_UNNATURAL\_MOTION} & Character or object motion is physically implausible. & Compare flow direction and magnitude with expected motion constraints. & \tblcode{constrain\_motion} \\
\tblcode{MO\_EVENT\_MISSING} & A required action or event does not occur in the expected segment. & Ask the VLM judge whether the specified action is visible. & \tblcode{inject\_action\_keyframe} \\
\tblcode{MO\_ACTION\_ORDER\_WRONG} & Actions occur in the wrong temporal order. & Compare detected action timestamps with the specification order. & \tblcode{reorder\_segments} \\
\tblcode{MO\_MOTION\_BLUR\_EXCESS} & Excessive motion blur degrades keyframes. & Use sharpness measures such as Laplacian variance. & \tblcode{reduce\_motion\_blur} \\
\tblcode{MO\_FROZEN\_FRAME} & The video is frozen when motion is expected. & Detect near-zero flow across the full segment. & \tblcode{regenerate\_segment} \\
\tblcode{MO\_SEGMENT\_BREAK} & A visual discontinuity occurs at a segment boundary. & Detect embedding-drift spikes at boundaries. & \tblcode{smooth\_segment\_transition} \\
\midrule
\multicolumn{4}{l}{Camera} \\
\tblcode{CA\_MOVE\_WRONG} & Camera movement deviates from the specified movement. & Compare optical-flow direction patterns with the expected movement type. & \tblcode{set\_camera\_movement} \\
\tblcode{CA\_SHOT\_TYPE\_WRONG} & The shot type does not match the specification. & Compare face or body box size ratios with the expected shot type. & \tblcode{adjust\_framing} \\
\tblcode{CA\_UNINTENDED\_ZOOM} & Unplanned zoom-in or zoom-out occurs. & Detect radial flow without a zoom specification. & \tblcode{remove\_zoom} \\
\tblcode{CA\_SHAKE} & Camera shake appears outside the specified movement. & Detect high-frequency oscillation in global optical flow. & \tblcode{stabilize\_camera} \\
\tblcode{CA\_WRONG\_ANGLE} & Camera angle deviates from the specification. & Ask the VLM judge to compare the observed angle with the target angle. & \tblcode{set\_camera\_angle} \\
\tblcode{CA\_ROLL} & Unintended camera roll appears. & Estimate the rotational component of optical flow. & \tblcode{remove\_camera\_roll} \\
\midrule
\multicolumn{4}{l}{Alignment} \\
\tblcode{AL\_PROP\_MISSING} & A required prop is not visible in the segment. & Use object detection or VLM verification for prop absence. & \tblcode{inject\_prop} \\
\tblcode{AL\_PROP\_WRONG\_PLACEMENT} & A required prop appears in the wrong region. & Compare the prop box with the expected screen region. & \tblcode{reposition\_prop} \\
\tblcode{AL\_TEXT\_PROMPT\_MISMATCH} & The generated content does not align with the prompt. & Measure text-image similarity and verify mismatches with a judge. & \tblcode{strengthen\_prompt\_guidance} \\
\tblcode{AL\_CHARACTER\_COUNT\_WRONG} & The number of visible characters does not match the specification. & Count faces or bodies and compare with the expected count. & \tblcode{adjust\_character\_count} \\
\tblcode{AL\_WRONG\_PROP} & A visible prop is not listed in the specification. & Detect high-confidence unlisted objects. & \tblcode{remove\_prop} \\
\tblcode{AL\_PROP\_OCCLUDED} & A required prop is present but substantially occluded. & Estimate visible area and occlusion of the prop box. & \tblcode{reposition\_prop} \\
\midrule
\multicolumn{4}{l}{Style} \\
\tblcode{ST\_COLOR\_SHIFT} & The color palette changes inconsistently across segments. & Measure HSV histogram distance between segments. & \tblcode{apply\_color\_grading} \\
\tblcode{ST\_ART\_STYLE\_INCONSISTENCY} & The art style shifts between segments. & Compare style embeddings across segments. & \tblcode{apply\_style\_transfer} \\
\tblcode{ST\_COMPRESSION\_ARTIFACT} & Visible compression artifacts degrade video quality. & Detect blocking, ringing, or quality-score degradation. & \tblcode{apply\_enhancement} \\
\tblcode{ST\_TEXTURE\_FLICKER} & Texture patterns flicker or shimmer unnaturally. & Measure high temporal frequency in texture regions. & \tblcode{smooth\_texture} \\
\tblcode{ST\_OVEREXPOSURE} & Significant frame regions are overexposed. & Count the fraction of saturated high-value pixels. & \tblcode{adjust\_exposure} \\
\tblcode{ST\_UNDEREXPOSURE} & Significant frame regions are underexposed. & Count the fraction of near-black pixels. & \tblcode{adjust\_exposure} \\
\end{longtable}
\endgroup


\section{Judge Protocol}
\label{app:judge}
The Stage-2 judge receives the top-$K$ candidate segments produced by Stage-1. For each
segment, the input includes the segment identifier, temporal bounds, keyframe paths, candidate
failure codes, and specification context, including required props, expected actions, camera
movement, shot type, and character anchors. The judge answers structured questions rather than
producing a free-form report. This design keeps the output comparable across local open-source
VLMs and hosted vision-capable models.

For identity failures, the judge checks whether the face, body appearance, clothing, or other
distinctive visual attributes of the main character change inconsistently across keyframes.
For alignment failures, it verifies whether required props are visible, whether they appear in
the expected screen region, and whether the generated content matches the prompt or structured
specification. For camera failures, it compares the observed camera movement, shot type, and
viewpoint angle with the requested control signal. For motion failures, it verifies whether
required actions occur, whether they occur in the specified order, and whether stutter, frozen
frames, or segment breaks are visible. For style and scene failures, it checks lighting, color,
compression artifacts, background consistency, and environment match.

The output is a structured object containing a natural-language answer for each question, a
confidence score, per-code probability updates, and a re-ranking of keyframes. The final
failure score is computed as
\[
\tilde{\sigma}_{s,c} = (1-\alpha)\sigma^{(1)}_{s,c} + \alpha\sigma^{(2)}_{s,c},
\]
where $\sigma^{(1)}_{s,c}$ is the Stage-1 score, $\sigma^{(2)}_{s,c}$ is the judge score,
and $\alpha=0.6$ by default. All judge calls record the model name, prompt, raw response,
token count when available, and output path to support reproducibility.


\section{Patch Template Examples}
\label{app:patch}
VideoGenDoctor compiles each diagnosis record into one or more patch actions. If a ShotIR
specification is available, the compiler emits field-level diffs. If the generator does not
expose a structured specification interface, the compiler emits generator-agnostic repair
plans that can be translated into prompt edits, reference-frame injection, segment-level
re-rendering, or post-processing steps.

\begin{verbatim}
{"action": "inject_reference_frame",
 "field": "character_id",
 "value": "<anchor_frame>",
 "reason": "ID_FACE_DRIFT detected at t=1.2s"}
\end{verbatim}

\begin{verbatim}
{"action": "set_camera_movement",
 "field": "camera_movement",
 "value": "<expected_camera_movement>",
 "reason": "CA_MOVE_WRONG detected in the localized span"}
\end{verbatim}

\begin{verbatim}
{"action": "inject_prop",
 "field": "props_required",
 "value": "<missing_prop>",
 "reason": "AL_PROP_MISSING verified by object or VLM evidence"}
\end{verbatim}

\begin{verbatim}
{"action": "regenerate_segment",
 "field": "segment",
 "value": "<t0,t1>",
 "reason": "MO_FROZEN_FRAME or MO_SEGMENT_BREAK is localized to this span"}
\end{verbatim}

Each patch specifies an action type, the affected target, an optional update value, and a short
rationale tied to the evidence span. When multiple records affect the same target, the compiler
groups them before re-rendering to reduce conflicts among patch actions.

\begin{table}[t]
\centering\scriptsize
\caption{Patch-to-generator adapter details. A repair plan is counted as adapter-executable only when it can be automatically translated into target-generator or editor inputs without manual rewriting. L0/L1 rows denote abstention or prompt-rewrite instructions rather than executable patches.}
\label{tab:patch_adapter_details}
\begin{adjustbox}{max width=\textwidth}
\begin{tabular}{p{0.18\textwidth}p{0.14\textwidth}p{0.10\textwidth}p{0.15\textwidth}p{0.24\textwidth}p{0.12\textwidth}}
\toprule
Repair-plan action & ShotIR field & Counted as & Adapter availability & Concrete generator/editor input & Fallback \\
\midrule
inject\_reference\_frame & character.identity\_anchor & L2 & L2 available for SVD; L1 partial for text-to-video models & reference image input + identity prompt strengthening & prompt-only repair plan \\
set\_camera\_movement & camera.movement & L1 & L1 partial & camera-control prompt rewrite; trajectory adapter when supported & no\_op\_uncertain if no camera control exists \\
set\_camera\_angle & camera.viewpoint & L1 & L1 partial & viewpoint prompt rewrite; optional camera-control token & repair suggestion only \\
inject\_prop & props\_required & L1 / L2 & L1 available; L2 partial with mask/reference input & prop phrase insertion + optional reference or mask & prompt edit \\
reposition\_prop & props.region & L1 & L1 partial & region-aware prompt rewrite or mask-conditioned edit & prompt-only repair plan \\
regenerate\_segment & temporal span & L2 / L3 & L2 available when segment rerendering is supported & segment-level rerender or temporal inpainting call & full-clip rerender \\
stabilize\_camera & camera.constraint & L1 / L2 & L1 partial & no-shake constraint or stabilization post-process & repair suggestion only \\
normalize\_frame\_rate & motion.continuity & L2 & L2 partial & frame interpolation or temporal smoothing post-process & regenerate\_segment \\
apply\_enhancement & style.quality & L2 & L2 available & enhancement / deblocking post-process & prompt edit for higher quality \\
fix\_background\_anchor & scene.background\_anchor & L1 & L1 partial & background-anchor prompt strengthening; optional reference image & prompt-only repair plan \\
no\_op\_uncertain & -- & L0 & L0 available & abstain from concrete repair & no action \\
\bottomrule
\end{tabular}
\end{adjustbox}
\end{table}

In the reported Human Pass@K evaluation, the concretely executed actions are the L2/L3 subset exposed by the active backend, primarily reference-frame injection, segment rerendering or temporal smoothing, enhancement or deblocking, and mask- or reference-conditioned prop insertion when supported. L0/L1 rows remain structured repair-plan instructions and are not counted as executable patches.

\section{Annotation Guide Summary}
\label{app:annotation}
Annotators review each perturbed video together with the original prompt or ShotIR
specification and the auto-assigned candidate failure codes. For each candidate, annotators
determine whether the failure is visually present, assign the tightest temporal span that contains
the visible deviation, and select representative keyframes when available. Annotators may also
add missing failures if the perturbation produces an additional visible artifact.

\begin{figure}[t]
\centering
\safeincludegraphics[width=\textwidth]{../figures/annotation_interface_appendix.png}
\caption{Annotation interface used for manual verification in VideoGenDoctor-Bench-v0. The tool exposes a sample navigator, side-by-side review of the original and perturbed videos, temporal evidence marking, structured failure verification cards, and free-form notes for adjudication.}
\label{fig:annotation_interface_appendix}
\end{figure}

The annotation record is stored as one JSON object per video and includes the video identifier,
annotator identifier, verified failure codes, confidence values, evidence spans, and free-form
notes. A failure is retained in the benchmark only when the annotator can verify it from the
rendered video and specification, without relying on perturbation metadata alone. Ambiguous
cases are handled conservatively: if a deviation is not visually identifiable, the candidate is
marked as not verified.

The experiment fixture includes a dual-annotated reliability subset of 120 videos and
768 candidate evidence spans. The subset contains both controlled perturbation samples and
naturally occurring real-failure samples. Larger multi-annotator validation with more than two
independent annotators remains future work.

\section{Example Diagnostic Report}
\label{app:diagnostic_report}

\begin{figure}[t]
\centering
\safeincludegraphics[width=\textwidth]{../figures/report.png}
\caption{Example VideoGenDoctor diagnostic report for a perturbed video. The report summarizes global scores, ranked failure records with confidence and temporal spans, representative keyframes, patch hints derived from the retained failure codes, and segment-level evidence used for localized review and repair.}
\label{fig:diagnostic_report}
\end{figure}
% End inlined from appendix.tex



\section{Direct VLM Baseline Prompt Templates}
\label{app:vlm_baselines}

The four direct VLM/LLM baselines differ only in their allowed information access and output schema. The free-form and structured VLM report baselines receive sampled frames and the specification. The VLM-to-patch baseline also receives sampled frames, the taxonomy names, and the repair-plan vocabulary, but not VideoGenDoctor's rule scores or patch compiler outputs. The Score+LLM-to-patch baseline receives scalar evaluator scores and the specification only; it does not receive sampled frames. These information boundaries are fixed before evaluation and are used for all samples.

\paragraph{F.1 Free-form VLM report prompt.}
The model receives the video specification, sampled keyframes with timestamps, and a short instruction: identify visible failures, explain the evidence, and propose repair suggestions. The response is then parsed into taxonomy codes and evidence spans for evaluation.

\paragraph{F.2 Structured VLM report prompt.}
The model receives the same input but must output JSON objects with fields \texttt{failure\_code}, \texttt{temporal\_span}, \texttt{affected\_target}, \texttt{confidence}, \texttt{evidence\_keyframes}, \texttt{rationale}, and \texttt{repair\_action}. Failures without visual evidence must be marked as uncertain rather than hallucinated.

\begin{verbatim}
You are evaluating whether a generated video follows the given specification.
For each visible failure, return a JSON record:
{
  "failure_code": "<taxonomy code>",
  "temporal_span": [t0, t1],
  "affected_target": "<character / prop / camera / motion / style / scene>",
  "confidence": <0-1>,
  "evidence_keyframes": ["frame@time", ...],
  "rationale": "short visual reason",
  "repair_action": "<patch action>"
}
Only report failures that are visually supported by the frames.
\end{verbatim}

\paragraph{F.3 VLM-to-patch prompt.}
The VLM-to-patch baseline receives the same video specification and sampled keyframes as the direct VLM diagnosis baselines, but it is asked to output repair actions directly rather than diagnosis records. This baseline is allowed to access the taxonomy names and the supported repair-plan vocabulary, but it is not given VideoGenDoctor's Stage-1 rule scores, candidate segments, or patch compiler outputs. It may infer temporal spans from the sampled frames, but it does not receive ground-truth spans or VideoGenDoctor-predicted spans. This setting tests whether a VLM can directly produce structured repair plans from the specification and visual evidence.

\begin{verbatim}
You are a video-generation repair planner.

Input:
1. A video specification or ShotIR-style instruction.
2. Sampled keyframes from the generated video, each with a timestamp.
3. The failure-code taxonomy names.
4. The supported repair-plan vocabulary.

Your task:
Inspect the sampled frames and the specification. Identify only visually supported
problems and output structured repair-plan actions. Do not write a free-form critique.
Do not invent failures that are not visible in the frames.

You may use the following failure-code groups:
- Identity: ID_FACE_DRIFT, ID_BODY_DRIFT, ID_WRONG_CHARACTER,
  ID_CLOTHING_CHANGE, ID_MULTIPLE_FACES, ID_NO_FACE_VISIBLE
- Camera: CA_MOVE_WRONG, CA_SHOT_TYPE_WRONG, CA_UNINTENDED_ZOOM,
  CA_SHAKE, CA_WRONG_ANGLE, CA_ROLL
- Motion: MO_JITTER, MO_FRAME_DROP, MO_UNNATURAL_MOTION,
  MO_EVENT_MISSING, MO_ACTION_ORDER_WRONG, MO_MOTION_BLUR_EXCESS,
  MO_FROZEN_FRAME, MO_SEGMENT_BREAK
- Alignment: AL_PROP_MISSING, AL_PROP_WRONG_PLACEMENT,
  AL_TEXT_PROMPT_MISMATCH, AL_CHARACTER_COUNT_WRONG,
  AL_WRONG_PROP, AL_PROP_OCCLUDED
- Style: ST_COLOR_SHIFT, ST_ART_STYLE_INCONSISTENCY,
  ST_COMPRESSION_ARTIFACT, ST_TEXTURE_FLICKER,
  ST_OVEREXPOSURE, ST_UNDEREXPOSURE
- Scene: SC_BG_INCONSISTENCY, SC_ENVIRONMENT_MISMATCH,
  SC_LIGHTING_SHIFT, SC_OBJECT_TELEPORT, SC_BG_FLICKER,
  SC_WRONG_TIME_OF_DAY

Supported repair-plan vocabulary:
- inject_reference_frame
- replace_character_anchor
- fix_character_consistency
- adjust_framing
- set_camera_movement
- set_camera_angle
- stabilize_camera
- remove_zoom
- remove_camera_roll
- inject_prop
- reposition_prop
- remove_prop
- adjust_character_count
- strengthen_prompt_guidance
- inject_action_keyframe
- reorder_segments
- regenerate_segment
- interpolate_frames
- normalize_frame_rate
- constrain_motion
- smooth_segment_transition
- reduce_motion_blur
- apply_color_grading
- apply_style_transfer
- apply_enhancement
- smooth_texture
- adjust_exposure
- fix_background_anchor
- set_environment
- normalize_lighting
- stabilize_scene_objects
- set_time_of_day
- no_op_uncertain

Return a JSON object with the following schema:
{
  "patches": [
    {
      "patch_id": "<unique patch id>",
      "failure_code": "<taxonomy code or uncertain>",
      "action": "<one item from the supported repair-plan vocabulary>",
      "target": "<character / prop / camera / motion / style / scene target>",
      "temporal_span": [<start_sec>, <end_sec>],
      "evidence_keyframes": ["frame@time", "..."],
      "confidence": <0-1>,
      "patch_fields": {
        "field": "<generator or ShotIR field to modify>",
        "value": "<new value, constraint, or reference to insert>"
      },
      "rationale": "<short visual reason>"
    }
  ],
  "global_notes": "<one-sentence summary>",
  "uncertain": <true or false>
}

Rules:
1. Output valid JSON only.
2. Every patch must use an action from the supported repair-plan vocabulary.
3. If no visually supported failure is visible, return:
   {"patches": [], "global_notes": "No visually supported repair action.", "uncertain": true}
4. If the temporal span is uncertain, estimate it from the nearest sampled frame
   timestamps and set "confidence" below 0.5.
5. Do not output patches for purely aesthetic preferences unless the specification
   explicitly requires them.
\end{verbatim}

\paragraph{F.4 Score+LLM-to-patch prompt.}
The Score+LLM-to-patch baseline receives scalar evaluator outputs and the original video specification, but it does not receive sampled frames. This baseline is intentionally weaker in visual grounding but represents a common production setting where low-level evaluators return scalar scores and an LLM converts low-scoring dimensions into repair suggestions. Because it cannot inspect frames, it is not allowed to claim frame-specific visual evidence. It may output approximate temporal spans only when the scalar evaluator provides segment-level scores; otherwise the temporal span must be marked as \texttt{null}. This makes the information boundary explicit and prevents unfair comparison with frame-based VLM baselines.

\begin{verbatim}
You are a repair planner for controllable video generation.

Input:
1. The original video specification or ShotIR-style instruction.
2. Scalar evaluator outputs, such as quality, alignment, motion, identity,
   camera, style, and scene-consistency scores.
3. Optional segment-level scalar scores if available.
4. The supported repair-plan vocabulary.

Important information boundary:
You do NOT receive sampled frames or the video itself. Therefore:
- Do not claim that a failure is visually observed.
- Do not mention specific visual evidence or keyframes.
- Only infer likely repair actions from low-scoring evaluator dimensions.
- Use "evidence_source": "scalar_score" for all patches.
- Use "temporal_span": null unless segment-level scores identify a specific interval.

Supported repair-plan vocabulary:
- inject_reference_frame
- replace_character_anchor
- fix_character_consistency
- adjust_framing
- set_camera_movement
- set_camera_angle
- stabilize_camera
- remove_zoom
- remove_camera_roll
- inject_prop
- reposition_prop
- remove_prop
- adjust_character_count
- strengthen_prompt_guidance
- inject_action_keyframe
- reorder_segments
- regenerate_segment
- interpolate_frames
- normalize_frame_rate
- constrain_motion
- smooth_segment_transition
- reduce_motion_blur
- apply_color_grading
- apply_style_transfer
- apply_enhancement
- smooth_texture
- adjust_exposure
- fix_background_anchor
- set_environment
- normalize_lighting
- stabilize_scene_objects
- set_time_of_day
- no_op_uncertain

Return a JSON object with the following schema:
{
  "patches": [
    {
      "patch_id": "<unique patch id>",
      "score_dimension": "<low-scoring evaluator dimension>",
      "inferred_failure_type": "<short inferred failure description>",
      "action": "<one item from the supported repair-plan vocabulary>",
      "target": "<character / prop / camera / motion / style / scene target>",
      "temporal_span": null,
      "evidence_source": "scalar_score",
      "triggering_scores": {
        "<score_name>": <score_value>
      },
      "confidence": <0-1>,
      "patch_fields": {
        "field": "<generator or ShotIR field to modify>",
        "value": "<new value, constraint, or reference to insert>"
      },
      "rationale": "<short reason based only on scalar scores and specification>"
    }
  ],
  "global_notes": "<one-sentence summary>",
  "uncertain": <true or false>
}

Rules:
1. Output valid JSON only.
2. Do not claim visual evidence because frames are not provided.
3. Do not output frame IDs or evidence keyframes.
4. Use temporal_span = null unless segment-level scalar scores are provided.
5. Every patch must use an action from the supported repair-plan vocabulary.
6. If the scalar scores do not indicate a clear repair target, return:
   {"patches": [], "global_notes": "No reliable score-based repair action.", "uncertain": true}
\end{verbatim}

\paragraph{F.5 Parsing and invalid-output handling.}
Invalid JSON outputs are repaired using a deterministic parser when possible. If the output cannot be parsed, the prediction is marked invalid. Failure codes outside the taxonomy are mapped to the nearest taxonomy code only when an exact synonym exists; otherwise they are counted as false positives. Predictions without temporal spans receive no localization credit. Hallucinated failures without visual evidence are counted as false positives. Repair actions outside the supported repair-plan vocabulary are marked invalid for repair-plan scoring and non-adapter-executable for adapter scoring.

\section{Blind Human Repair Evaluation Protocol}
\label{app:human_repair}
Raters are shown the specification and one output video, without method names or diagnosis records. They answer whether the video satisfies the specification overall, whether the target failure is resolved, whether the repair introduces new artifacts, and how useful the patch is on a 1--5 scale. Human Pass@$K$ is the fraction of samples judged successful after at most $K$ repair attempts.

\begin{table}[t]
\centering
\caption{Blind human repair evaluation rubric for patch usefulness.}
\label{tab:human_rubric}
\begin{tabular}{cl}
\toprule
Score & Criterion \\
\midrule
1 & The patch does not address the target failure. \\
2 & The patch is related to the failure but is not actionable or is largely insufficient. \\
3 & The patch is actionable but incomplete or requires substantial manual interpretation. \\
4 & The patch directly addresses the target failure with minor ambiguity. \\
5 & The patch directly resolves the target failure and preserves surrounding valid content. \\
\bottomrule
\end{tabular}
\end{table}

A repaired video is marked as pass if the target failure is resolved and no severe new artifact is introduced. Human Pass@K is the fraction of samples judged successful after at most $K$ repair attempts. New artifacts are marked when raters observe a visible degradation introduced by repair that was not present in the original generated video. Each repaired video is rated by 3 independent raters. We report the majority-vote pass/fail label, the mean patch-usefulness score, and the majority-vote new-artifact label.

\begin{table}[t]
\centering\scriptsize
\caption{Error analysis on ambiguous and real-normal samples. The table focuses on no-op behavior and wrong-target repair planning. Best values are boldfaced for interpretable comparison columns; the no\_op\_uncertain rate is reported descriptively rather than bold-ranked.}
\label{tab:no_op_wrong_target_analysis}
\begin{adjustbox}{max width=\textwidth}
\begin{tabular}{lccccc}
\toprule
Method & no\_op\_uncertain rate & Correct no-op rate & Missed necessary repair & Wrong-target patch & Over-confident patch \\
\midrule
Score-only & 0.286 & 0.524 & 0.119 & 0.171 & 0.321 \\
VLM structured report & 0.219 & 0.476 & \textbf{0.096} & 0.229 & 0.396 \\
VLM temporal proposal + patch & 0.257 & 0.518 & 0.108 & 0.200 & 0.350 \\
VideoGenDoctor-diagnosis & 0.552 & \textbf{0.738} & 0.142 & 0.095 & 0.181 \\
VideoGenDoctor-full & 0.514 & 0.724 & 0.156 & \textbf{0.081} & \textbf{0.164} \\
\bottomrule
\end{tabular}
\end{adjustbox}
\end{table}

\begin{table}[t]
\centering\scriptsize
\caption{Qualitative case-study summary. Cases are sampled from controlled, real-failure, and ambiguous stress subsets.}
\label{tab:case_study_summary}
\begin{tabular}{lccc}
\toprule
Case type & \#Cases & Majority-judged successful / appropriate & Typical outcome \\
\midrule
Successful structured repair plan & 24 & 18 & localized repair plan resolves target failure \\
Ambiguous case with abstention & 24 & 16 & no\_op\_uncertain avoids over-repair \\
Over-repair / wrong-target failure & 24 & 7 & concrete patch targets wrong object or span \\
\bottomrule
\end{tabular}
\end{table}

\section{Bootstrap Confidence Intervals}
\label{app:bootstrap}
We use 10,000 video-level bootstrap resamples and report percentile 95\% confidence intervals. For the controlled fixture, resampling is performed over 324 videos. For the real-failure subset, resampling is performed over 240 videos. For blind human repair evaluation, resampling is performed over the stratified human-rated subset. For paired comparisons, we resample videos with replacement and compute the distribution of metric differences between two methods. Small point-estimate differences are not treated as meaningful unless the paired confidence interval excludes zero.

\begin{table}[t]
\centering\scriptsize
\caption{Leave-one-code-family-out rule ablation. For the held-out family, dedicated family-specific rules are disabled; only generic embedding drift, optical-flow anomaly, scene-change peaks, and structured VLM verification are retained.}
\label{tab:leave_one_code_family}
\begin{adjustbox}{max width=\textwidth}
\begin{tabular}{lcccc}
\toprule
Held-out family & Macro-F1 & tIoU@0.5 & Top-3 hit & Main degradation source \\
\midrule
Identity & 0.782 & 0.579 & 0.536 & face/body-specific anchors removed \\
Camera & 0.754 & 0.546 & 0.501 & camera-flow templates disabled \\
Motion & 0.721 & 0.512 & 0.468 & jitter / segment-break rules disabled \\
Alignment & 0.801 & 0.603 & 0.558 & prop detector and placement rules disabled \\
Style & 0.836 & 0.641 & 0.590 & compression/style detectors disabled \\
Scene & 0.748 & 0.528 & 0.487 & background-anchor rules disabled \\
\midrule
Average & 0.774 & 0.568 & 0.523 & -- \\
\bottomrule
\end{tabular}
\end{adjustbox}
\end{table}

The large degradation under leave-one-code-family-out ablation confirms that VideoGenDoctor should not be interpreted as an unconstrained open-ended failure discovery system. The experiment clarifies which portions of the performance depend on family-specific detectors and which portions remain supported by generic evidence signals and VLM verification.

\section{Real-Failure Construction Manifest}
\label{app:real_failure_manifest}

\begin{table}[t]
\centering\scriptsize
\caption{Generator configuration metadata for the real-failure and real-normal subsets. Hosted or version-changing endpoints are logged with access date and endpoint identifier.}
\label{tab:generator_config}
\begin{adjustbox}{max width=\textwidth}
\begin{tabular}{lccccccccc}
\toprule
Generator & Checkpoint / endpoint & fps & Frames & Resolution & Steps & Guidance & Scheduler & Seed range & Reference source / prompt rule \\
\midrule
CogVideoX & CogVideoX1.5-5B & 16 & 128 & 768$\times$432 & 50 & 6.0 & DPMSolver & 10000--10059 & ShotIR-to-prompt template v1 \\
Wan & Wan2.1-T2V-14B & 16 & 128 & 768$\times$432 & 50 & 5.0 & FlowMatchEuler & 18000--18059 & ShotIR-to-prompt template v1 \\
Stable Video Diffusion & SVD-XT-1.1 & 14 & 112 & 768$\times$432 & 40 & 7.5 & EulerDiscrete & 26000--26059 & reference frame from ShotIR anchor \\
HunyuanVideo & HunyuanVideo-1.5 & 16 & 128 & 768$\times$432 & 50 & 7.0 & FlowMatchEuler & 34000--34059 & multi-shot prompt template v1 \\
\bottomrule
\end{tabular}
\end{adjustbox}
\end{table}

When a generator does not expose one of these configuration fields, the corresponding manifest value is set to \texttt{null} rather than omitted. For hosted or version-changing endpoints, we log the endpoint identifier, access date, raw request payload, and raw response metadata.

To make the real-failure subset auditable and reproducible, each generated video is accompanied by a JSON manifest. The manifest records the generator identity, checkpoint, input specification, generation configuration, output location, screening decision, verified failure codes, temporal evidence spans, and annotator identifiers. This protocol turns the real-failure subset from an informal collection of generated examples into a traceable evaluation resource.

Each manifest entry contains the following fields: \texttt{generator}, \texttt{checkpoint}, \texttt{prompt\_id}, \texttt{input\_type}, \texttt{shotir\_spec}, \texttt{shotir\_to\_prompt\_converted\_text}, \texttt{reference\_frame\_id}, \texttt{seed}, \texttt{resolution}, \texttt{fps}, \texttt{num\_frames}, \texttt{sampling\_steps}, \texttt{guidance\_scale}, \texttt{output\_path}, \texttt{screening\_status}, \texttt{verified\_codes}, \texttt{evidence\_spans}, and \texttt{annotator\_ids}. The \texttt{screening\_status} field takes one of four values: \texttt{generated}, \texttt{retained\_candidate}, \texttt{verified}, or \texttt{excluded\_ambiguous}. Only entries marked as \texttt{verified} are included in the reported real-failure evaluation.

\begin{verbatim}
{
  "video_id": "real_wan_00037",
  "generator": "Wan",
  "checkpoint": "Wan2.1-T2V-14B",
  "prompt_id": "shotir_prompt_037",
  "input_type": "text / ShotIR-to-prompt",
  "shotir_spec": {
    "shot_id": "S037",
    "duration_sec": 8.0,
    "characters": [
      {
        "id": "person_A",
        "description": "young man wearing a blue hoodie",
        "identity_anchor": "ref_person_A_004"
      }
    ],
    "required_props": [
      {
        "id": "red_water_bottle",
        "description": "red water bottle held in the right hand",
        "expected_presence": "visible throughout the shot"
      }
    ],
    "camera": {
      "shot_type": "medium shot",
      "movement": "slow left-to-right tracking",
      "viewpoint": "front three-quarter view"
    },
    "action_order": [
      "person_A enters the park path",
      "person_A jogs while holding the red water bottle",
      "person_A slows down near the bench"
    ],
    "style": {
      "lighting": "natural daylight",
      "visual_style": "realistic video"
    },
    "scene": {
      "location": "urban park path",
      "background_anchor": "trees and bench remain stable"
    }
  },
  "shotir_to_prompt_converted_text": "A realistic 8-second medium shot of a young man wearing a blue hoodie jogging on an urban park path while holding a red water bottle in his right hand. The camera performs a slow left-to-right tracking movement from a front three-quarter view. The trees and bench in the background should remain stable. Natural daylight. The man enters the path, jogs with the red water bottle, and slows down near the bench.",
  "reference_frame_id": null,
  "seed": 18473,
  "resolution": "768x432",
  "fps": 16,
  "num_frames": 128,
  "sampling_steps": 50,
  "guidance_scale": 7.5,
  "output_path": "real_failures/wan/real_wan_00037.mp4",
  "screening_status": "verified",
  "verified_codes": [
    "AL_PROP_MISSING",
    "CA_MOVE_WRONG",
    "MO_SEGMENT_BREAK"
  ],
  "evidence_spans": [
    {
      "failure_code": "AL_PROP_MISSING",
      "target": "red_water_bottle",
      "start_sec": 3.25,
      "end_sec": 6.75,
      "evidence_keyframes": [
        "frame_052@3.25s",
        "frame_080@5.00s",
        "frame_108@6.75s"
      ],
      "annotator_confidence": 0.91,
      "notes": "The specified red water bottle is not visible after the character begins jogging."
    },
    {
      "failure_code": "CA_MOVE_WRONG",
      "target": "camera_movement",
      "start_sec": 0.00,
      "end_sec": 4.00,
      "evidence_keyframes": [
        "frame_000@0.00s",
        "frame_032@2.00s",
        "frame_064@4.00s"
      ],
      "annotator_confidence": 0.84,
      "notes": "The camera remains mostly static rather than performing the requested left-to-right tracking movement."
    },
    {
      "failure_code": "MO_SEGMENT_BREAK",
      "target": "temporal_continuity",
      "start_sec": 5.50,
      "end_sec": 5.94,
      "evidence_keyframes": [
        "frame_088@5.50s",
        "frame_095@5.94s"
      ],
      "annotator_confidence": 0.78,
      "notes": "A visible discontinuity occurs near the transition to the final action."
    }
  ],
  "annotator_ids": [
    "ann_02",
    "ann_05"
  ]
}
\end{verbatim}

For privacy and release safety, file paths may be anonymized, but the manifest must preserve all fields required to reproduce the generation configuration, reconstruct the ShotIR-to-prompt conversion, and audit the verified evidence spans. When a generator does not expose a configuration field such as \texttt{guidance\_scale}, the field is set to \texttt{null} rather than omitted.

\newpage
\input{checklist.tex}

\end{document}
